\documentclass[11pt,a4paper]{article}

% Packages
\usepackage[utf8]{inputenc}
\usepackage[T1]{fontenc}
\usepackage{amsmath,amssymb,amsfonts}
\usepackage{graphicx}
\graphicspath{{images/lemon/}}
\usepackage{booktabs}
\usepackage{array}
\usepackage{geometry}
\usepackage{setspace}
\newcommand{\degree}{$^\circ$}
\usepackage{textcomp}
\usepackage{threeparttable}
\usepackage{hyperref}
\usepackage{xcolor}
\usepackage{float}
\usepackage{caption}
\usepackage{tcolorbox}
\usepackage{longtable}

% Unicode character support
\usepackage{newunicodechar}
\newunicodechar{⁵}{\textsuperscript{5}}
\newunicodechar{⁴}{\textsuperscript{4}}
\newunicodechar{⁰}{\textsuperscript{0}}
\newunicodechar{ⁿ}{\textsuperscript{n}}
\newunicodechar{≤}{\leq}
\newunicodechar{≈}{\approx}
\newunicodechar{−}{\ensuremath{-}}
\newunicodechar{₀}{$_0$}
\newunicodechar{₁}{$_1$}
\newunicodechar{₂}{$_2$}
\newunicodechar{₃}{$_3$}
\newunicodechar{₄}{$_4$}
\newunicodechar{₅}{$_5$}
\newunicodechar{₆}{$_6$}
\newunicodechar{₇}{$_7$}
\newunicodechar{₈}{$_8$}
\newunicodechar{₉}{$_9$}
\newunicodechar{⁷}{\textsuperscript{7}}
\newunicodechar{√}{\sqrt{}}
\newunicodechar{ⱼ}{\textsubscript{j}}
\newunicodechar{ₖ}{_k}
\newunicodechar{≥}{\geq{}}
\newunicodechar{⁶}{\textsuperscript{6}}
\newunicodechar{⁻}{\textsuperscript{-}}

% Language support
\usepackage[english]{babel}

% Page setup
\geometry{margin=1in}
\onehalfspacing

% Hyperref setup
\hypersetup{
    colorlinks=true,
    linkcolor=blue,
    citecolor=blue,
    urlcolor=blue
}

% Custom commands
\newcommand{\phisym}{\varphi}
\newcommand{\fzero}{f_0}

% Colored boxes for key equations and insights
\newtcolorbox{keyequation}{
    colback=blue!5!white,
    colframe=blue!75!black,
    fonttitle=\bfseries,
    title=Key Equation
}

\newtcolorbox{keyinsight}{
    colback=green!5!white,
    colframe=green!50!black,
    fonttitle=\bfseries,
    title=Key Insight
}

\newtcolorbox{criticalpoint}{
    colback=red!5!white,
    colframe=red!50!black,
    fonttitle=\bfseries,
    title=Critical Methodological Point
}

% Bibliography
\usepackage[style=apa,backend=biber,sorting=nyt]{biblatex}
\addbibresource{lemon_refs.bib}

\title{Dominant Peaks, Fragile Metrics: Separating Architectural Signal from Methodological Artifact in Phi-Lattice EEG Analysis}

\author{Michael Lacy\\
\small Independent Researcher\\
\small \href{mailto:michael.lacy@gmail.com}{michael.lacy@gmail.com}}

\date{\today}

\begin{document}
\maketitle

% =====================================================================
\begin{abstract}
% =====================================================================

The golden ratio ($\phisym = 1.618$) has been proposed to organize EEG spectral peaks into a logarithmic lattice. Across three datasets---LEMON ($N = 202$), the Dortmund Vital Study ($N = 608$), and EEGMMIDB ($N = 109$)---we find that dominant oscillatory peaks sit closer to $\phisym$-lattice positions than chance, with per-subject mean lattice distance converging to $\bar{d} = 0.069$ (range: 0.0689--0.0693) across all three datasets. Comparing nine logarithmic bases using coverage-normalized degree-3 noble positions under per-octave aperiodic modeling (overlap-trim extraction), $\phisym$ ranks 1st in all three datasets ($\bar{d}/\bar{d}_{\text{null}} = 0.901$--$0.917$, all $p < 10^{-3}$), with alignment concentrated at three positions in aggregate cross-band analysis (boundary, attractor, noble$_1$); higher-degree positions carry no individual signal. This alignment is age-invariant ($r = -0.004$, $p = 0.92$; Dortmund ages 20--70 continuous), fatigue-invariant ($0.914 \to 0.909$ across a 2-hour battery), and cognitively silent in the dataset with available cognitive measures (LEMON: 0/100+ tests survive FDR; $N = 202$, 80\% power to $|r| = 0.20$). Critically, $\phisym$'s rank-1 status requires per-octave aperiodic modeling; under standard whole-range extraction, $\phisym$ ranks 7th of 9 ($p = 0.83$). Replication with method-independent aperiodic separation (IRASA, eBOSC) is the highest-priority confirmation target.

A cross-dataset sign flip resolves the $\fzero$ vs IAF/2 ambiguity: in Dortmund, $\fzero = 7.83$~Hz and IAF/2 $\approx 5.0$~Hz lie in opposite directions from theta. Under eyes-closed conditions, theta converges on $\fzero$ ($d = -0.385$, $p < 10^{-19}$) while diverging from IAF/2 ($d = +0.351$, $p < 10^{-16}$), excluding all IAF-derived models.

However, five diagnostic tests show aggregate enrichment metrics are dominated by extraction parameters: structural scores collapse under FOOOF range changes and lose significance at higher peak counts. A five-year longitudinal arm ($N = 208$) confirms the lattice describes population-level architecture, not individual fingerprints: group $\bar{d}$ is invariant, but within-subject ICC is negative ($-0.25$ to $-0.36$). The alignment predicts neither cognitive performance (tested in LEMON only) nor organization beyond the four strongest oscillations---but $\fzero$ is an absorbing state in a potential landscape confirmed by the two-dataset sign flip.

\end{abstract}

\bigskip
\noindent\textbf{Significance Statement.}
The golden ratio has been proposed to organize brain rhythms, but prior evidence relied on aggregate spectral metrics vulnerable to analysis choices. We show that the single strongest oscillation in each canonical frequency band falls closer to logarithmic lattice positions than chance for bases near the golden ratio---an effect that replicates across three datasets totaling 919 subjects (ages 20--77; EEGMMIDB from our prior work \parencite{Lacy2026paper2}, LEMON as primary analysis, Dortmund as out-of-sample replication), is age-invariant, fatigue-invariant, and nearly doubles under eyes-closed conditions ($d = 0.75$). The golden ratio ($\phisym$) ranks 1st among nine candidate bases in all three datasets under coverage-normalized degree-3 comparison, with the second-place base varying across datasets (1.7, 3/2, $\sqrt{2}$). A cross-dataset sign flip definitively identifies the convergence target as the Schumann resonance fundamental ($f_0 = 7.83$~Hz) rather than an alpha subharmonic. A five-year longitudinal follow-up ($N = 208$) reveals that population-level alignment is invariant across sessions while individual alignment shows no test-retest reliability (ICC~$< 0$), identifying this as a species-level architectural constant rather than an individual trait. Yet this regularity predicts nothing about cognition in the one dataset with available cognitive measures (LEMON, $N = 202$; over 100 tests, zero FDR survivors across three independent metrics). The $\phisym$-specificity claim is conditional on per-octave aperiodic modeling: under standard whole-range FOOOF, $\phisym$ does not separate from competing bases. Confirmation with method-independent aperiodic separation (IRASA, eBOSC) is the critical next step. The brain's frequency architecture carries a logarithmic-lattice signature anchored to a geophysical frequency, but that signature is a structural fossil rather than a functional lever.

\bigskip
\noindent\textbf{Keywords:} golden ratio, EEG spectral peaks, phi-lattice, FOOOF/specparam, frequency architecture, dominant oscillations, eyes-closed EEG, cross-base comparison, LEMON dataset, Dortmund Vital Study, EEGMMIDB, Schumann resonance, cross-dataset replication

\newpage
\tableofcontents
\newpage

% =====================================================================
\section{Introduction}
% =====================================================================

\subsection{Logarithmic frequency organization in neural oscillations}

Classical EEG frequency bands---delta through gamma---span roughly logarithmic frequency ranges, with each band covering approximately one octave \parencite{Buzsaki2004, Penttonen2003}. This logarithmic scaling is ubiquitous in sensory systems (Weber-Fechner law, cochlear tonotopy, visual spatial frequency tuning; \cite{Buzsaki2006}) and has led several authors to propose that the golden ratio specifically organizes the spacing between EEG frequency bands \parencite{Roopun2008, Pletzer2010, Klimesch2012, Kramer2022}.

The $\phisym$-lattice framework formalizes this proposal. Given an anchor frequency $\fzero$ and the golden ratio $\phisym = (1+\sqrt{5})/2$, any frequency $f$ maps to a lattice coordinate:
\begin{keyequation}
\begin{equation}
u = \log_\phisym(f / \fzero) \bmod 1
\end{equation}
\end{keyequation}
\noindent Four positions within each $\phisym$-octave carry special dynamical significance, derived from $\phisym$'s continued-fraction expansion ($\phisym = [1; 1, 1, 1, \ldots]$):
\begin{itemize}
    \item \textbf{Boundary} ($u = 0.000$): the $\phisym$-octave edge, where adjacent octaves meet
    \item \textbf{Noble$_2$} ($u = 0.382$): the lower noble position ($1/\phisym^2$)
    \item \textbf{Attractor} ($u = 0.500$): the geometric midpoint of the $\phisym$-octave
    \item \textbf{Noble$_1$} ($u = 0.618$): the upper noble position ($1/\phisym$)
\end{itemize}
These positions predict that EEG spectral peaks should cluster at specific frequencies determined by $\fzero$ and $\phisym$, rather than being distributed uniformly in log-frequency space. The equidistribution of these positions follows from the three-distance theorem applied to Fibonacci fractions \parencite{Slater1967}.

\subsection{The Schumann anchor hypothesis and prior claims}

The lattice requires an anchor frequency $\fzero$. The Schumann resonance fundamental ($\sim$7.83~Hz; \cite{Nickolaenko2014}) has been proposed as a natural anchor, placing the lattice fundamental at the theta-alpha boundary. Prior work on the EEGMMIDB dataset ($N = 109$; \cite{Lacy2026paper2}; data from \cite{Schalk2004, Goldberger2000}) found the global optimum of a joint ($\fzero$, base) grid search at $\fzero = 8.5$~Hz, base $= \phisym$, with a structural score SS $= 45.6$ and bootstrap confidence intervals excluding the next-best base. The same study demonstrated that irrational bases outperform rational bases ($d = 1.2$--$1.6$) and identified a sharp sigmoid transition at $\fzero \sim 7.5$~Hz below which $\phisym$'s advantage disappears.

A key vulnerability remained untested: the EEGMMIDB analysis used 1.86 million FOOOF-detected peaks across 109 subjects, creating potential sensitivity to FOOOF extraction parameters. This vulnerability is not dataset-specific: re-extracting EEGMMIDB peaks at [1,\,75]~Hz collapses SS from $+45.6$ to $-23.6$, the same pattern we observe in LEMON (Section~\ref{sec:freq_range}). The fragility of aggregate enrichment metrics is a methodology problem, not a replication failure. We additionally test the dominant-peak findings in the Dortmund Vital Study ($N = 608$, 64 channels, ages 20--70; \cite{GudiMindermann2024}), providing a third independent dataset with continuous age coverage and a pre/post-task design that enables fatigue-invariance testing.

\subsection{This paper: a progressive diagnostic}

Rather than testing a single hypothesis, we apply five complementary diagnostic approaches---each designed to probe a different potential confound---followed by a sixth analysis targeting the most fundamental claim: whether dominant oscillatory generators sit at lattice positions.

\begin{enumerate}
    \item \textbf{Standard enrichment replication}: Does the all-peaks structural score replicate in LEMON?
    \item \textbf{Amplitude-weighted enrichment}: Does weighting by peak prominence rescue sensitivity to \texttt{max\_n\_peaks}?
    \item \textbf{FOOOF frequency range and fitting window}: Are enrichment metrics stable across extraction parameters?
    \item \textbf{$\fzero$ parameter sensitivity}: Is the age effect lattice erosion or IAF positioning?
    \item \textbf{Model-free pairwise ratios}: Does $\phisym$-specificity survive without an anchor frequency?
    \item \textbf{Dominant-peak alignment}: Does the strongest oscillation per band sit at a lattice position?
\end{enumerate}

Each analysis is motivated by the previous result. The narrative moves from apparent confirmation toward increasingly honest evaluation---and ultimately reveals what actually survives. The dominant-peak findings are then replicated in a third dataset, the Dortmund Vital Study ($N = 608$, ages 20--70 continuous), providing definitive cross-dataset confirmation across 919 subjects and resolving key ambiguities---particularly the $\fzero$ vs IAF/2 disambiguation---that the LEMON data alone could not cleanly separate.

% =====================================================================
\section{Methods}
% =====================================================================

\subsection{Dataset: LEMON}

The Leipzig Study for Mind-Body-Emotion Interactions (LEMON; \cite{Babayan2019}) provides $N = 202$ healthy adults: 136 young (20--35 years) and 66 elderly (59--77 years), with an intentional age gap at 36--58 years. EEG was recorded with a 62-channel BrainVision system at 2500~Hz. Preprocessing included ICA artifact rejection, average re-reference, and bandpass filtering at 0.5--45~Hz. Each participant completed approximately 8 minutes each of eyes-open (EO) and eyes-closed (EC) resting-state recording.

Cognitive measures included the CVLT (verbal learning), TMT-A/B (processing speed, set shifting), TAP (alertness, working memory, inhibition), LPS (fluid reasoning), WST (crystallized intelligence), and RWT (verbal fluency). TMT and TAP-Alert were log-transformed for normality; nine tests entered all cognitive analyses. Of the 202 subjects, $N = 182$ entered the all-peaks diagnostic analysis (Part~I); all 202 were used for dominant-peak analyses (Part~II). The difference reflects 20 subjects held out for internal validation of the diagnostic pipeline. For eyes-closed analyses, $N = 203$ (one additional subject had usable EC but not EO data). Under standard whole-range FOOOF extraction, some subjects lack a valid dominant peak in one or more bands, reducing the available sample: $N = 196$ for standard EO and $N = 179$ for standard EC (Table~\ref{tab:fourway}). The overlap-trim method recovers all subjects ($N = 202$ EO, $N = 203$ EC) because per-octave aperiodic fitting is more robust. These sample-size variations are noted in each analysis.

\subsection{Dataset: Dortmund Vital Study}

The Dortmund Vital Study provides $N = 608$ healthy adults aged 20--70 years with continuous age coverage (no gap). EEG was recorded with a 64-channel BrainVision system at 1000~Hz, downsampled to 250~Hz with 50~Hz notch filtering. Preprocessing included ICA artifact rejection and average re-reference. Each participant completed four resting-state blocks: eyes-closed pre-task (EC-pre), eyes-open pre-task (EO-pre), eyes-closed post-task (EC-post), and eyes-open post-task (EO-post), where pre/post refer to a 2-hour cognitive battery. This design enables within-subject tests of both condition sensitivity (EC vs EO) and fatigue invariance (pre vs post). A longitudinal arm provides $N = 208$ subjects who returned approximately five years later for an identical protocol (session~2), enabling direct assessment of within-subject stability (Section~\ref{sec:longitudinal}). No cognitive test scores were available for the present analysis. Data are publicly available on OpenNeuro (ds005385; \cite{GudiMindermann2024}).

\subsection{Dataset: EEGMMIDB}

The EEG Motor Movement/Imagery Dataset (EEGMMIDB; $N = 109$ subjects, 64 channels, 160~Hz) serves as a cross-reference for population-level alignment. EEGMMIDB was the primary analysis dataset in our prior work \parencite{Lacy2026paper2}, where the $\phisym$-lattice framework was initially developed; its inclusion here is therefore a within-framework consistency check rather than a blind replication. FOOOF extraction parameters and analysis methods follow the same protocol as LEMON; full dataset description appears in \textcite{Lacy2026paper2}. Data from \textcite{Schalk2004} and \textcite{Goldberger2000}.

\subsection{Peak extraction variants}

Five FOOOF (SpecParam; \cite{Donoghue2020}) extraction configurations were used across the diagnostic cascade (Table~\ref{tab:extraction}). All used per-channel Welch PSDs \parencite{Welch1967} (2-second windows, 50\% overlap) with SpecParam's default settings except where noted. For methodological considerations regarding oscillation detection, see \textcite{Donoghue2022methods}.

\begin{table}[H]
\centering
\caption{FOOOF peak extraction configurations.}
\label{tab:extraction}
\begin{tabular}{lllcl}
\toprule
Label & Freq range & Fitting window & \texttt{max\_n\_peaks} & Purpose \\
\midrule
Standard & [1,\,45] Hz & Global & 20 & Baseline replication \\
Extended & [1,\,85] Hz & Global & 20 & Frequency range sensitivity \\
Per-$\phisym$-octave & [1,\,45] Hz & Per $\phisym$-octave & 20 & Window placement \\
Overlap-trim & [1,\,45] Hz & Target $+$ 0.5 $\phisym$-octave & 20 & Edge artifact control \\
max40 & [1,\,45] Hz & Global & 40 & Peak count sensitivity \\
\bottomrule
\end{tabular}
\end{table}

\subsection{All-peaks metrics}

\textbf{Structural Score (SS).} For all extracted peaks, we compute $u = \log_\phisym(f/\fzero) \bmod 1$, then estimate kernel density enrichment at each of the four lattice positions using a Gaussian kernel ($\sigma = 0.03$), following \textcite{Lacy2026paper1}. Enrichment is defined as the percentage deviation from the uniform null expectation. The structural score is:
\begin{equation}
\text{SS} = -E_{\text{boundary}} + E_{\text{attractor}} + \frac{1}{2}(E_{\text{noble}_1} + E_{\text{noble}_2})
\end{equation}
where positive SS indicates peak concentration at attractors and nobles with depletion at boundaries.

\textbf{Amplitude-weighted SS.} Replaces the uniform kernel contribution of each peak with a weight proportional to FOOOF peak power. Four weight transforms were tested: rank ($w_i = \text{rank}(p_i)/N$), z-score (with negative values clamped to zero), linear ($10^{p_i}$), and raw.

\textbf{Within-band shuffle null.} Lattice coordinates are permuted uniformly within $\phisym$-octave bands (preserving band identity; a gamma peak stays in gamma), over 1{,}000 iterations. The resulting $z$-score indexes whether within-band peak organization exceeds chance.

\subsection{Dominant-peak metric}

For each subject and each canonical band (delta: 1--4, theta: 4--8, alpha: 8--13, gamma: 30--45~Hz):
\begin{enumerate}
    \item Select the highest-power FOOOF peak within the band across all channels
    \item Compute the lattice coordinate $u = \log_\phisym(f/\fzero) \bmod 1$
    \item Compute distance to nearest lattice position: $d = \min_{p \in \mathcal{P}} \text{circ}(u, p)$, where $\mathcal{P}$ is the position set (degree 2: $\{0, 0.382, 0.5, 0.618\}$; degree 3: $\{0, 0.236, 0.382, 0.5, 0.618, 0.764\}$; see Section~\ref{sec:crossbase_methods})
    \item Subject-level alignment: $\bar{d} = \frac{1}{4}\sum_{\text{bands}} d_{\text{band}}$
\end{enumerate}
where $\text{circ}(a, b) = \min(|a - b|, 1 - |a - b|)$ is the circular distance on $[0, 1)$. The individual-level alignment test (Section~\ref{sec:dominant_peak}) uses degree-2 positions; the cross-base comparison (Section~\ref{sec:crossbase}) uses degree~3.

Beta (13--30~Hz) is excluded because it spans the second harmonic of alpha ($\sim$20~Hz), making the dominant beta peak harmonically coupled to alpha via the sensorimotor mu rhythm \parencite{Pfurtscheller1999} rather than reflecting an independent oscillatory generator. This coupling is confirmed empirically: (i)~beta and alpha all-peaks age correlations are nearly identical ($r = -0.170$ vs $-0.168$, $\Delta r = 0.002$; Table~\ref{tab:band_decomp}), consistent with both tracking individual alpha frequency (IAF) decline rather than independent aging trajectories; and (ii)~including a harmonically dependent band would violate the independence assumption of the $\bar{d} = \frac{1}{N}\sum d_{\text{band}}$ metric, effectively double-weighting the alpha contribution.

Pre-specified: $\fzero = 7.83$~Hz (Schumann fundamental). No optimization.

\subsubsection{Cross-base comparison: higher-degree noble positions and coverage normalization}
\label{sec:crossbase_methods}

To test whether dominant-peak alignment is specific to base $\phisym$ or a property of any logarithmic lattice, we compare nine exponential bases $b \in \{1.4, \sqrt{2}, 3/2, \phisym, 1.7, 1.8, 2, e, \pi\}$. Each base generates its own set of noble positions at $(1/b)^k$ and $1 - (1/b)^k$ for integer degree $k = 1, 2, \ldots, k_{\max}$, plus the universal boundary ($u = 0$) and attractor ($u = 0.5$). Candidate positions within $0.02$ of each other or of $0/1$ are deduplicated (keeping the lower-degree representative).

At the standard degree $k_{\max} = 2$, $\phisym$'s self-similarity ($1/\phisym = \phisym - 1$, hence $(1/\phisym)^2 = 1 - 1/\phisym$) causes its degree-2 nobles to collapse onto degree-1, yielding only 4 distinct positions vs 5--6 for other bases. This creates a position-count asymmetry that inflates $\bar{d}$ for $\phisym$. We resolve this by extending to $k_{\max} = 3$, where $\phisym$ gains genuinely new positions at $(1/\phisym)^3 = 0.236$ and $1 - (1/\phisym)^3 = 0.764$, yielding 6 distinct positions.

Different bases still have different position counts (6--8) and spacings at degree 3, so raw $\bar{d}$ retains a coverage confound. We normalize by computing:
\begin{equation}
\bar{d}_{\text{norm}} = \bar{d} \;/\; \bar{d}_{\text{null}}
\end{equation}
where $\bar{d}_{\text{null}}$ is the analytical expected distance to the nearest position for a uniformly distributed point on $[0, 1)$ given the base's specific position layout. For $N_p$ positions at locations $\{p_1, \ldots, p_{N_p}\}$ on the unit circle, $\bar{d}_{\text{null}} = \frac{1}{4} \sum_{i} g_i^2$, where $g_i$ is the gap between the $i$-th and $(i+1)$-th positions (circular, so $g_{N_p}$ wraps from $p_{N_p}$ to $p_1$). Values $\bar{d}_{\text{norm}} < 1$ indicate peaks closer to positions than chance; the ratio controls for both position count and spacing.

Significance is assessed by a one-sided paired $t$-test of per-subject $\bar{d}$ against $\bar{d}_{\text{null}}$, separately for each base.

\subsection{Pairwise ratio analysis (anchor-free)}

For all peak pairs within each subject (same channel), we compute pairwise frequency ratios and map each to $u = \log_{\text{base}}(\text{ratio}) \bmod 1$. We test clustering at lattice positions for nine bases: $\{1.4, \sqrt{2}, 3/2, \phisym, 1.7, 1.8, 2, e, \pi\}$. This removes $\fzero$ entirely---the ratio structure is intrinsic.

\subsection{Statistical framework}

Age effects were assessed by Pearson and Spearman correlations with continuous age and Cohen's $d$ for young vs elderly group comparisons. Cognitive prediction used partial correlations controlling for age, with Benjamini-Hochberg FDR correction \parencite{Benjamini1995} across all tests. Null models included uniform within-band, Gaussian (matched parameters), cross-band shuffle, and within-band shuffle (1{,}000 permutations). Cross-condition stability was evaluated with paired tests and position transition matrices.

\textbf{Sign convention.} All Cohen's $d$ values are reported as unsigned effect sizes (alignment strength) except where directional interpretation is essential (Section~\ref{sec:disambiguation}), where negative $d$ indicates convergence toward the target and positive $d$ indicates divergence from it.

% =====================================================================
\section{Results}
% =====================================================================

\subsection*{Part I: The Diagnostic---All-Peaks Metrics Under Scrutiny}

\subsection{Standard enrichment partially replicates}

At $\fzero = 7.83$~Hz with standard [1,\,45]~Hz extraction ($N = 182$ analysis subjects):
\begin{itemize}
    \item SS $= +11.7$ (positive, weaker than EEGMMIDB's $+45.6$)
    \item Attractor enrichment: $+4.1\%$
    \item Boundary enrichment: $-6.4\%$ (expected depletion, but $p = 0.390$)
    \item Noble enrichment: $+9.7\%$ ($p = 0.332$)
    \item Split-half reliability: ICC $= 0.804$
\end{itemize}

\begin{keyinsight}
The sole statistically significant finding: attractor enrichment correlates with age ($r = -0.195$, $p = 0.008$, $d = 0.41$). All other cognitive and demographic tests are null after FDR correction. The remainder of the diagnostic probes whether this age effect is genuine.
\end{keyinsight}

\subsection{Amplitude weighting does not rescue the signal}

The \texttt{max\_n\_peaks}~$= 40$ setting eliminates the attractor age trend ($r = -0.110$, $p = 0.138$), suggesting weak FOOOF peaks act as noise. If prominent peaks carry the signal, amplitude weighting should restore it. Instead, all four weight transforms (rank, z-score, linear, raw) \textit{attenuate} the age correlation rather than strengthening it. Prominence mediation is null: peak power does not mediate the age $\to$ compliance pathway (bootstrap indirect effect CI includes zero). Weighting noise louder does not help because the signal was never in peak amplitudes.

\subsection{FOOOF extraction parameters dominate enrichment}
\label{sec:diagnostic_cascade}

\subsubsection{Frequency range sensitivity}
\label{sec:freq_range}

Re-extracting peaks at [1,\,85]~Hz to cover the full gamma $\phisym$-octave produces a sign reversal. (The [1,\,85]~Hz FOOOF range exceeds the nominal 45~Hz preprocessing lowpass; because this filter rolls off gradually rather than as a brick wall, FOOOF detects peaks in attenuated spectral regions above 45~Hz, as confirmed by the raw vs preprocessed comparison in Appendix~B, where preprocessed data yields peaks up to 71~Hz.)
\begin{itemize}
    \item SS collapses: $+11.7 \to -67.4$
    \item EEGMMIDB shows the same: SS $= +45.6$ at [1,\,50] $\to -23.6$ at [1,\,75]
    \item The Dortmund Vital Study (Section~\ref{sec:measurement}) shows the same pattern under standard extraction: delta anti-alignment under EC ($d = 0.31$), confirming the 1/$f$ knee mechanism is dataset-general
\end{itemize}

\begin{criticalpoint}
The mechanism is FOOOF's fixed peak budget \parencite{Ostlund2022}. The \texttt{max\_n\_peaks}~$= 20$ parameter is a hard ceiling per channel. When the fitting range extends above 45~Hz, newly visible high-frequency peaks compete for the same 20 slots, displacing peaks below 45~Hz. The lattice coordinate distribution changes not because the EEG changed, but because the algorithm allocated its budget differently. At LEMON's spectral resolution (0.122~Hz), 54\% of channels already saturate the 20-peak ceiling at [1,\,45]~Hz; extending the range exacerbates the problem.
\end{criticalpoint}

An intermediate approach---fitting FOOOF separately within each $\phisym$-octave band---produces even more extreme artifacts (SS $= +119.3$ at primary offset, $-84.9$ at 0.5 $\phisym$-octave offset; see Appendix~\ref{app:peroctave}).

\subsubsection{Overlap-trim greatly reduces but does not eliminate artifacts}

The overlap-trim method (fit on wider windows, keep only interior peaks) improves stability:
\begin{itemize}
    \item Primary: SS $= +12.9$; offset: SS $= +5.8$ ($\Delta = 7.1$ vs 204.2 for per-$\phisym$-octave)
    \item Per-subject correlation primary $\leftrightarrow$ offset: $r = 0.72$--$0.80$ (vs near-zero for per-$\phisym$-octave)
    \item But attractor enrichment still flips sign: $+3.0\% \to -7.7\%$
\end{itemize}
Age correlations are stable across both offsets (SS $\sim$ age: primary $r = +0.374$, offset $r = +0.423$, both $p < 0.001$), but the direction is positive---elderly show \textit{higher} compliance, reversed from the global [1,\,45] result.

\subsection{The age effect is IAF positioning, not lattice erosion}

\subsubsection{Band decomposition}

Decomposing the attractor age trend by frequency band reveals its source (Table~\ref{tab:band_decomp}).

\begin{table}[H]
\centering
\caption{Band decomposition of attractor age correlations.}
\label{tab:band_decomp}
\begin{tabular}{lcccc}
\toprule
Band & Attractor $r(\text{age})$ & $p$ & Shuffle-$z$ $r(\text{age})$ & $p$ \\
\midrule
Alpha & $-0.168$ & .024 & $-0.198$ & .007 \\
Beta & $-0.170$ & .022 & $-0.191$ & .010 \\
Delta+Theta & $-0.124$ & .101 & $-0.136$ & .079 \\
Gamma & $+0.155$ & .037 & $-0.010$ & .893 \\
\bottomrule
\end{tabular}
\end{table}

The delta+theta null is decisive: approximately 90\% of the attractor age effect originates in the alpha and beta bands. Beta is harmonically coupled to alpha (not independent evidence). The effect is IAF slowing \parencite{Scally2018, Klimesch1999} expressed in $\phisym$-lattice coordinates.

\subsubsection{Per-subject $\fzero$ correction}

\begin{table}[H]
\centering
\caption{Attractor age correlation under different anchor frequencies.}
\label{tab:f0_correction}
\begin{tabular}{lccc}
\toprule
$\fzero$ setting & $r(\text{age}, E_{\text{attractor}})$ & $p$ & Interpretation \\
\midrule
$8.5$ (EEGMMIDB optimum) & $-0.223$ & .002 & Fixed anchor \\
$7.83$ (Schumann) & $-0.195$ & .008 & Fixed anchor \\
IAF-derived (per-subject) & $\mathbf{+0.167}$ & .025 & \textbf{Reversed} \\
$\fzero^*$ (optimized) & $+0.092$ & .215 & Null \\
\bottomrule
\end{tabular}
\end{table}

The mechanism: young adults (mean IAF $= 10.15$~Hz) map to $u = 0.365$ (near noble$_2$); elderly adults (mean IAF $= 9.96$~Hz) map to $u = 0.325$ (drifting toward boundary). Per-subject $\fzero$ correction rotates out the IAF projection, which \textit{is} the effect.

\subsubsection{$\fzero$ shift confirmation}

\begin{keyinsight}
Shifting $\fzero$ from 7.83 to 8.50~Hz with overlap-trim extraction eliminates the age effect:

\medskip
\begin{tabular}{lcccc}
\toprule
$\fzero$ & SS $\sim$ age $r$ & $p$ & Attractor $\sim$ age $r$ & Cohen's $d$ \\
\midrule
7.83 & $+0.374$ & $<.001$ & $+0.314$ & $+0.79$ \\
8.50 & $\mathbf{+0.022}$ & $\mathbf{.754}$ & $\mathbf{-0.007}$ & $\mathbf{+0.001}$ \\
\bottomrule
\end{tabular}

\medskip
\noindent Confirmed across all 8 conditions (2 $\fzero$ $\times$ 2 data types $\times$ 2 offsets). The age effect is a geometric accident of where $\fzero$ places IAF in lattice coordinates.
\end{keyinsight}

\subsection{Pairwise ratio analysis: $\phisym$ ranks 4th}

When the anchor frequency is removed entirely:
\begin{itemize}
    \item RWT (verbal fluency) $\sim$ $z_{\text{alpha}}$ correlation appears at 7 of 9 bases; base 3/2 beats $\phisym$ ($r = 0.309$ vs $0.227$)
    \item Five bases survive FDR: 3/2, $\sqrt{2}$, $\phisym$, 1.4, 1.7---the cognitive effect is generic alpha-band spectral organization
    \item $z_{\text{global}} \sim$ age: $\phisym$ ($r = -0.259$) and $e$ ($r = -0.253$) strongest, both IAF-mediated
    \item Within-elderly cognitive results all fail FDR (0/36), many with wrong signs
\end{itemize}

\begin{criticalpoint}
\textbf{Summary of Part~I.} Every all-peaks metric examined is dominated by extraction parameters, IAF projection, or both. The age correlation is $\fzero$ placement; the structural score is FOOOF range; the enrichment is fitting window edges. $\phisym$ ranks 4th when the anchor is removed. The diagnostic has eliminated every claim based on aggregate spectral peak counts.
\end{criticalpoint}

Part~I establishes that aggregate spectral-peak metrics are methodological artifacts: they change sign with extraction parameters and rank $\phisym$ no higher than 4th among nine bases. This motivates a fundamentally different approach. Instead of counting hundreds of FOOOF-detected peaks---most of which reflect weak spectral fluctuations---Part~II restricts attention to one dominant peak per canonical band: the four strongest oscillatory generators per brain, which survive regardless of extraction parameters. The question becomes: do these four peaks sit at $\phisym$-lattice positions?

% ── Part II ──────────────────────────────────────────────────────────

\subsection*{Part II: What Survives---Dominant-Peak Architecture}

\subsection{Dominant-peak alignment: individual-level architecture}
\label{sec:dominant_peak}

At pre-specified $\fzero = 7.83$~Hz, using overlap-trim extraction (Section~2.2), individual subjects have dominant peaks closer to $\phisym$-lattice positions than expected by chance.

Per-subject $\bar{d}$: mean $= 0.069$, SD $= 0.028$, vs analytical null expectation $= 0.080$.

\begin{table}[H]
\centering
\caption{Individual-level alignment tests.}
\begin{tabular}{lcl}
\toprule
Test & Statistic & $p$-value \\
\midrule
One-sample $t$ & $t = -5.68$ & $p < 10^{-7}$ \\
Wilcoxon signed-rank & --- & $p < 10^{-7}$ \\
Cohen's $d$ & $0.40$ & --- \\
\bottomrule
\end{tabular}
\end{table}

This is not a population-averaging artifact. The effect is present in individual brains.

\textbf{Three-dataset convergence.} The Dortmund Vital Study ($N = 608$) replicates this finding under eyes-closed conditions: per-subject $\bar{d} = 0.069$ (EC-pre), Cohen's $d = 0.391$ ($p < 10^{-20}$)---identical to LEMON EO to the third decimal place. The EEGMMIDB dataset ($N = 109$) yields $\bar{d} = 0.069$, $d = 0.377$. Across three independent datasets totaling 919 subjects, recorded with different equipment, at different laboratories, with different preprocessing pipelines, the per-subject mean lattice distance converges to $\bar{d} = 0.069$ (range: 0.0689--0.0693).

\textbf{Per-band decomposition.} Theta drives the strongest alignment (Mann-Whitney $p < 10^{-6}$), followed by alpha ($p = 0.005$). Delta and gamma are not individually significant by the one-sample $t$-test used for theta and alpha, but show significant departures from uniformity by KS test against a simulated null distribution (all $p < 10^{-9}$), indicating distributional clustering at lattice positions that does not reach the per-subject effect size of the theta and alpha bands.

\textbf{Fraction individually significant.} 13 of 202 subjects (6.4\%) reach $p < 0.05$ against a per-subject bootstrap null, vs 10.1 expected at the nominal 5\% rate. The binomial test is non-significant ($p = 0.30$)---the effect is a population-level shift in alignment, not a subpopulation of ``$\phisym$-aligned'' individuals.

\subsection{$\phisym$-specificity through noble positions}
\label{sec:crossbase}

\textbf{The position-count problem.} Each base $b$ generates noble positions at $(1/b)^k$ and $1 - (1/b)^k$ for integer degree $k$. At degree $k \leq 2$ (the standard noble set), most bases obtain 5--6 distinct positions (boundary, attractor, and 2--4 nobles). However, $\phisym$ is unique: its self-similarity property $1/\phisym = \phisym - 1$ causes degree-2 nobles to collapse onto degree-1. Specifically, $(1/\phisym)^2 = 1 - 1/\phisym$ and $1 - (1/\phisym)^2 = 1/\phisym$, so the six candidate positions reduce to four: $\{0, 0.382, 0.5, 0.618\}$. This is a mathematical consequence of $\phisym$'s continued-fraction expansion $[1; 1, 1, 1, \ldots]$ and means that any cross-base comparison using degree-2 nobles is asymmetric: $\phisym$ has fewer positions to ``catch'' nearby peaks, inflating its mean distance $\bar{d}$.

\textbf{Degree-2 comparison: $\phisym$ ranks last.} Using each base's natural degree-$\leq 2$ positions, $\phisym$ ranks 9th of 9 bases in per-subject $\bar{d}$ in both LEMON ($N = 202$) and EEGMMIDB ($N = 109$). Bases with 6 positions (e.g., 1.7, 1.8) achieve lower $\bar{d}$ simply because their positions tile the unit interval more densely. This result is uninformative about $\phisym$-specificity.

\textbf{Higher-degree nobles resolve the asymmetry.} At degree $k = 3$, $\phisym$ gains genuinely new positions: $(1/\phisym)^3 = 0.236$ and $1 - (1/\phisym)^3 = 0.764$, yielding six distinct positions $\{0, 0.236, 0.382, 0.5, 0.618, 0.764\}$. These higher-degree nobles arise from the same continued-fraction structure as the degree-1 nobles and reflect deeper mode-locking avoidance in $\phisym$-scaled systems. With degree-3 positions, $\phisym$ jumps from rank 9 to rank 3 (raw $\bar{d}$) in all three datasets (Table~\ref{tab:crossbase}). Note that $\bar{d}$ drops from 0.069 (Section~\ref{sec:dominant_peak}, degree 2, 4 positions) to 0.042 (degree 3, 6 positions) because more positions mechanically reduce the distance to the nearest one for all bases; the cross-base ranking, not the absolute $\bar{d}$, is the informative comparison.

\textbf{Coverage normalization: $\phisym$ ranks 1st.} Different bases still have different numbers of degree-3 positions (6--8), and their geometric spacing differs, so raw $\bar{d}$ retains a coverage confound. We normalize by computing $\bar{d}/\bar{d}_{\text{null}}$, where $\bar{d}_{\text{null}}$ is the expected mean distance for uniformly distributed points given each base's specific position layout. This ratio measures how much closer peaks fall to lattice positions than chance, controlling for both position count and spacing. Under this normalization, $\phisym$ ranks 1st in all three datasets (Table~\ref{tab:crossbase}), with the second-place base varying (1.7 in LEMON, 3/2 in Dortmund, $\sqrt{2}$ in EEGMMIDB).

\begin{table}[H]
\centering
\caption{Cross-base dominant-peak alignment using degree-3 symmetric noble positions across three datasets. Each base generates positions at $(1/b)^k$ and $1 - (1/b)^k$ for $k = 1, 2, 3$, plus boundary ($u = 0$) and attractor ($u = 0.5$), with deduplication at 0.02 minimum separation. $N_{\text{pos}}$: distinct positions. $\bar{d}/\bar{d}_{\text{null}}$: coverage-normalized ratio (lower = closer to lattice positions than chance). Sorted by average rank across datasets. $\phisym$ ranks 1st in all three datasets; the second-place base varies (1.7, 3/2, $\sqrt{2}$).}
\label{tab:crossbase}
\begin{tabular}{lccccccc}
\toprule
& & \multicolumn{2}{c}{LEMON ($N = 202$)} & \multicolumn{2}{c}{Dortmund ($N = 608$)} & \multicolumn{2}{c}{EEGMMIDB ($N = 109$)} \\
\cmidrule(lr){3-4} \cmidrule(lr){5-6} \cmidrule(lr){7-8}
Base & $N_{\text{pos}}$ & $\bar{d}/\bar{d}_{\text{null}}$ & Rank & $\bar{d}/\bar{d}_{\text{null}}$ & Rank & $\bar{d}/\bar{d}_{\text{null}}$ & Rank \\
\midrule
$\phisym$ & 6 & \textbf{0.917} & \textbf{1} & \textbf{0.914} & \textbf{1} & \textbf{0.901} & \textbf{1} \\
3/2 & 8 & 0.977 & 5 & 0.919 & 2 & 0.944 & 3 \\
1.7 & 8 & 0.919 & 2 & 0.951 & 4 & 0.992 & 5 \\
$\sqrt{2}$ & 6 & 0.985 & 6 & 0.950 & 3 & 0.940 & 2 \\
2 & 6 & 0.951 & 3 & 1.008 & 7 & 0.982 & 4 \\
1.8 & 8 & 0.958 & 4 & 0.990 & 5 & 1.015 & 7 \\
$\pi$ & 8 & 1.039 & 7 & 0.995 & 6 & 1.003 & 6 \\
1.4 & 6 & 1.040 & 8 & 1.083 & 9 & 1.136 & 8 \\
$e$ & 8 & 1.073 & 9 & 1.079 & 8 & 1.185 & 9 \\
\bottomrule
\end{tabular}
\end{table}

$\phisym$ ranks 1st in all three datasets: LEMON ($\bar{d}/\bar{d}_{\text{null}} = 0.917$, $p < 10^{-4}$), Dortmund EC-pre ($0.914$, $p < 10^{-11}$), and EEGMMIDB ($0.901$, $p < 10^{-3}$), meaning dominant peaks fall 8--10\% closer to $\phisym$-lattice positions than expected by chance (Figure~\ref{fig:three_dataset}). The second-place base varies across datasets: 1.7 in LEMON ($0.919$), 3/2 in Dortmund ($0.919$), and $\sqrt{2}$ in EEGMMIDB ($0.940$). Among the six-position bases ($\phisym$, $\sqrt{2}$, 2, 1.4), $\phisym$ significantly outperforms all others in all datasets ($p < 0.001$), ruling out a generic coverage explanation among equal-count bases.

\textbf{Why $\phisym$'s positions work.} The six degree-3 $\phisym$ positions $\{0, 0.236, 0.382, 0.5, 0.618, 0.764\}$ have the lowest null $\bar{d}_{\text{null}} = 0.0455$ among all six-position bases, meaning they tile the unit interval most uniformly. This is not coincidental: $\phisym$'s noble positions are consecutive Fibonacci fractions, which are optimally equidistributed by the three-distance theorem \parencite{Slater1967}. Bases with 8 positions at degree 3 (e.g., 1.8, 1.7, 3/2) have mechanically lower raw $\bar{d}$ due to denser coverage, but after normalization $\phisym$ still outperforms them in all three datasets. Several eight-position bases show significant alignment ($p < 0.05$) in individual datasets, indicating that the alignment signal is not exclusively $\phisym$-specific---but $\phisym$ is the only base that ranks 1st consistently across all three samples.

\textbf{Position-count decomposition.}
\label{sec:poscount_caveat}
The signal attenuates monotonically with position count (Figure~\ref{fig:poscount}). At degree 2 (4 positions), Cohen's $d \approx 0.38$--$0.40$ across all three datasets. At degree 3 (6 positions), $d \approx 0.28$--$0.34$. At degree 4 (8 positions), $d$ drops to $\sim$0.07 (LEMON, Dortmund) or $\sim$0 (EEGMMIDB). To distinguish genuine signal dilution from the coarse-partition effect (any structured distribution shows larger departures at coarser partitions), we test whether individual positions carry significant alignment (Table~\ref{tab:per_position}).

\begin{table}[H]
\centering
\caption{Per-position enrichment across three datasets. Enrichment (\%) is the deviation of KDE density at each position from the uniform null ($10{,}000$ permutations, Gaussian kernel $\sigma = 0.03$). Positions with $p < 0.05$ are bolded. The degree-3 positions ($u = 0.236$, $u = 0.764$) are consistently \textit{depleted}, not enriched, meaning peaks actively avoid these lattice coordinates. Note: values are aggregate cross-band figures pooling peaks from all frequency bands; individual bands may show distinct position preferences.}
\label{tab:per_position}
\begin{tabular}{llrrrrrrr}
\toprule
& & \multicolumn{2}{c}{LEMON ($N = 202$)} & \multicolumn{2}{c}{Dortmund ($N = 608$)} & \multicolumn{2}{c}{EEGMMIDB ($N = 109$)} \\
\cmidrule(lr){3-4} \cmidrule(lr){5-6} \cmidrule(lr){7-8}
Position & $u$ & Enrich. & $z$ & Enrich. & $z$ & Enrich. & $z$ \\
\midrule
\textbf{Boundary} & 0.000 & $\mathbf{+45\%}$ & $\mathbf{4.4}$ & $\mathbf{+42\%}$ & $\mathbf{7.1}$ & $+0\%$ & $0.0$ \\
Noble$_3$ & 0.236 & $-22\%$ & $-2.2$ & $-20\%$ & $-3.4$ & $-27\%$ & $-1.9$ \\
Noble$_2$ & 0.382 & $-6\%$ & $-0.5$ & $-8\%$ & $-1.4$ & $-27\%$ & $-1.9$ \\
\textbf{Attractor} & 0.500 & $\mathbf{+21\%}$ & $\mathbf{2.0}$ & $\mathbf{+26\%}$ & $\mathbf{4.4}$ & $\mathbf{+40\%}$ & $\mathbf{2.9}$ \\
\textbf{Noble$_1$} & 0.618 & $\mathbf{+37\%}$ & $\mathbf{3.7}$ & $\mathbf{+54\%}$ & $\mathbf{9.3}$ & $\mathbf{+92\%}$ & $\mathbf{6.5}$ \\
Inv-Noble$_3$ & 0.764 & $-25\%$ & $-2.4$ & $-45\%$ & $-7.6$ & $-26\%$ & $-1.9$ \\
\bottomrule
\end{tabular}
\end{table}

In the aggregate cross-band analysis, three positions carry consistent signal across all three datasets: noble$_1$ ($u = 0.618$, the strongest), attractor ($u = 0.500$), and boundary ($u = 0.000$; significant in LEMON and Dortmund but null in EEGMMIDB). Individual frequency bands may exhibit distinct position preferences (e.g., theta boundary enrichment driven by $\fzero$ convergence; Table~\ref{tab:position_stability}). Noble$_2$ ($u = 0.382$) is consistently null or depleted ($-6\%$ to $-27\%$ across datasets), despite being one of the four core degree-2 positions from the original $\phisym$-lattice framework. This three-of-four asymmetry means the effective signal-carrying lattice is $\{0, 0.5, 0.618\}$---boundary, attractor, and the upper noble---rather than the full degree-2 set. Whether this reflects a genuine three-position architecture or a power limitation (noble$_2$ lying closer to attractor than noble$_1$ does, potentially suffering competitive suppression from the adjacent enriched position) remains an open question. Critically, the degree-3 positions are \textit{anti-aligned}: per-subject distance to the two degree-3 positions alone is significantly \textit{greater} than chance ($d = -0.38$ to $-0.49$, all $p < 10^{-5}$; Table~\ref{tab:deg3_critical}). Peaks actively avoid these lattice coordinates.

\begin{table}[H]
\centering
\caption{Per-subject alignment to degree-3 positions only ($u = 0.236$, $u = 0.764$). Negative Cohen's $d$ indicates peaks are farther from these positions than expected under uniformity.}
\label{tab:deg3_critical}
\begin{tabular}{lccccc}
\toprule
Dataset & $\bar{d}_{\text{obs}}$ & $\bar{d}_{\text{null}}$ & Cohen's $d$ & $t$ & $p$ \\
\midrule
LEMON & 0.139 & 0.125 & $-0.38$ & $5.42$ & $1.7 \times 10^{-7}$ \\
EEGMMIDB & 0.141 & 0.125 & $-0.41$ & $4.30$ & $3.7 \times 10^{-5}$ \\
Dortmund & 0.143 & 0.125 & $-0.49$ & $12.16$ & $1.3 \times 10^{-30}$ \\
\bottomrule
\end{tabular}
\end{table}

This resolves the interpretive ambiguity of the position-count decomposition. The monotonic attenuation from $d \approx 0.40$ (degree 2) to $d \approx 0.28$ (degree 3) to $d \approx 0.07$ (degree 4) reflects the coarse-partition effect, not signal dilution: the degree-3 and degree-4 positions carry no signal (and in fact carry anti-signal), but including them mechanically reduces $\bar{d}$ because they are additional ``catching'' positions that happen to lie near some peaks by chance. In the aggregate analysis, the alignment signal is concentrated at three positions: noble$_1$, attractor, and boundary. Noble$_2$ is null; higher-degree nobles are repellers. Per-band enrichment patterns may differ from this aggregate ordering, as the diagnostic in Part~I demonstrates for other aggregate metrics.

One exception: noble$_3$ ($u = 0.236$) shows delta-specific enrichment in 2/3 datasets (Dortmund: $+34\%$, $z = 2.9$, $p = 0.003$; EEGMMIDB: $+94\%$, $z = 3.4$, $p = 0.002$; LEMON: $+28\%$, $z = 1.3$, $p = 0.10$), consistent with the delta dominant peak falling near $f_0 \times \phisym^{-3} \approx 1.84$~Hz. This delta-specific signal is overwhelmed by depletion in the other three bands.

Degree 3 remains the principled choice for cross-base comparison---the highest degree at which $\phisym$ gains genuinely non-collapsing positions---but the per-position evidence indicates that $\phisym$'s rank-1 advantage at degree 3 derives from the coverage-normalization properties of its position spacing \parencite{Slater1967}, not from signal at the degree-3 positions themselves.

The depletion has a concrete spectral explanation. In theta's $\phisym$-octave ($\fzero/\phisym = 4.84$~Hz to $\fzero = 7.83$~Hz), noble$_3$ ($u = 0.236$) maps to $\sim$5.4~Hz---the spectral gap between the $\fzero/\phisym$ subharmonic and the region where theta peaks actually cluster ($\sim$6--7.5~Hz). No dominant oscillation inhabits this gap, so theta depletes the position. The same logic applies band by band: each degree-3 coordinate falls in the spectral no-man's-land between adjacent band generators, where only delta (the weakest and least spatially consistent band) shows any enrichment. The three signal-carrying positions succeed precisely because each has a band whose dominant generator clusters there (theta$\to$boundary, alpha$\to$attractor/noble$_1$, gamma$\to$noble$_1$; Section~\ref{sec:ec_fourband}); the degree-3 positions fail because no band fills their spectral gap.

\textbf{Extraction sensitivity.} The dominant-peak metric uses one peak per conventional frequency band, eliminating sensitivity to $\phisym$-octave partitions and \texttt{max\_n\_peaks} settings. However, the degree-3 ranking is extraction-dependent. Under overlap-trim extraction, $\phisym$ ranks 1st ($\bar{d}/\bar{d}_{\text{null}} = 0.917$, $p < 10^{-4}$); under standard whole-range extraction, $\phisym$ ranks 7th of 9 ($\bar{d}/\bar{d}_{\text{null}} = 1.020$, $p = 0.83$---\textit{not significant}). The individual-level alignment test (degree-2) shows the same pattern: $d = 0.40$ ($p < 10^{-7}$) under overlap-trim vs $d = 0.12$ ($p = 0.08$) under standard extraction. Under standard extraction, bases 1.8, 1.7, and 2 rank 1st--3rd ($\bar{d}/\bar{d}_{\text{null}} = 0.82$--$0.84$), meaning they benefit more from the extraction change than $\phisym$ does.

The difference is driven by delta and gamma bands, where per-octave aperiodic modeling better isolates genuine oscillatory peaks from the 1/$f$ background (Section~3.11). As noted in Section~4.2, this is the methodologically recommended approach when the aperiodic spectrum deviates from a single power law \parencite{Donoghue2020}. All dominant-peak results in this paper use overlap-trim extraction; the cross-dataset replication in EEGMMIDB and Dortmund (also overlap-trim) confirms that the rank-1 result is not LEMON-specific, but it is extraction-specific.

\subsection{Alignment is age-invariant and cognitively silent}
\label{sec:cognitive}

\textbf{Age.} No age effect on dominant-peak alignment:

\begin{table}[H]
\centering
\begin{tabular}{lcl}
\toprule
Test & Statistic & $p$ \\
\midrule
$\bar{d} \sim$ age (Pearson) & $r = +0.054$ & 0.45 \\
$\bar{d} \sim$ age (Spearman) & $\rho = +0.077$ & 0.27 \\
Young vs elderly ($t$-test) & $d = 0.12$ & 0.42 \\
\bottomrule
\end{tabular}
\end{table}

\noindent The lattice describes a fixed architectural property that does not degrade with aging---a dissociation from the age-related changes in aperiodic exponent documented by \textcite{Voytek2015}. This contrasts sharply with the all-peaks attractor erosion ($r = -0.195$), which Section~3.4 showed was IAF positioning.

\textbf{Dortmund confirms age-invariance with continuous coverage.} The LEMON result is limited by its age gap (36--58 years absent). The Dortmund Vital Study provides continuous coverage from 20--70 years ($N = 601$ after excluding 7 subjects with missing age metadata): $r(\bar{d}, \text{age}) = -0.004$ ($p = 0.92$), young ($<$35) vs old ($>$55) Cohen's $d = 0.008$ ($p = 0.99$). Age is definitively unrelated to lattice alignment.

\textbf{Cognition.} Zero-order, age-partialed, per-band, and within-elderly correlations between $\bar{d}$ and nine cognitive measures all fail FDR correction (Table~\ref{tab:cognitive_null}).

\begin{table}[H]
\centering
\caption{Cognitive prediction across three independent metrics.}
\label{tab:cognitive_null}
\begin{tabular}{lccc}
\toprule
Test family & $N$ tests & FDR survivors & Largest $|r|$ \\
\midrule
Zero-order ($\bar{d}$ vs cog) & 9 & 0 & WST $r = 0.090$ \\
Age-partialed ($\bar{d}$ vs cog) & 9 & 0 & WST $r = 0.086$ \\
Per-band ($d_{\text{band}}$ vs cog) & 36 & 0 & $\delta$--WST $r = 0.218$, $p_{\text{FDR}} = 0.066$ \\
Within-elderly ($\bar{d}$ vs cog) & 9 & 0 & CVLT $r = 0.230$ \\
\bottomrule
\end{tabular}
\end{table}

This is the third independent metric yielding a cognitive null within LEMON, after all-peaks enrichment and within-band $z$-scores. No metric predicts cognitive performance in the one dataset with available cognitive measures.

\textbf{Power analysis.} With $N = 202$, the minimum detectable correlation at 80\% power ($\alpha = 0.05$, two-tailed) is $|r| = 0.196$. We can rule out moderate-or-larger cognitive associations; only small effects ($|r| < 0.20$) remain compatible with the data. The LEMON cognitive battery spans verbal learning (CVLT), processing speed and set-shifting (TMT-A/B), alertness, working memory, and inhibition (TAP), fluid reasoning (LPS), crystallized intelligence (WST), and verbal fluency (RWT)---nine measures covering the major cognitive domains. The null is not a consequence of narrow measurement.

\subsection{Measurement sensitivity vs architectural invariance}
\label{sec:measurement}

The Dortmund dataset's four within-subject conditions (EC-pre, EC-post, EO-post, EO-pre) create a natural experiment for measurement sensitivity: the same subjects, recorded with the same equipment on the same day, produce different spectral signal-to-noise ratios across conditions. Two complementary findings emerge (Figure~\ref{fig:measurement}).

\textbf{Degree-2 is measurement-sensitive.} Within Dortmund, individual-level alignment (Cohen's $d$) varies monotonically with condition: EC-pre $d = 0.391$, EC-post $d = 0.265$, EO-post $d = 0.213$, EO-pre $d = 0.037$ (not significant). This gradient tracks alpha peak prominence---EC conditions produce sharper spectral peaks and correspondingly stronger individual-level alignment. The EO-pre null ($d = 0.037$, $p = 0.18$) is not evidence of absent architecture; it reflects insufficient spectral SNR to detect individual-level alignment in short EO recordings. This explains why detecting dominant-peak alignment requires either EC conditions or per-octave extraction: both increase the precision with which dominant peaks are isolated.

\textbf{Degree-3 is architecturally invariant.} In contrast, the coverage-normalized metric $\bar{d}/\bar{d}_{\text{null}}$ for $\phisym$ at degree 3 is relatively stable across conditions: Dortmund EC-pre $= 0.914$, EC-post $= 0.909$, EO-post $= 0.955$, EO-pre $= 0.983$; LEMON EO $= 0.917$; EEGMMIDB $= 0.901$. $\phisym$ ranks 1st in all conditions except EO-pre (rank 2, behind 1.8 by 0.002). The EC values cluster near 0.91 and EO values near 0.97, consistent with EC conditions providing better spectral SNR. This dissociation is informative: the architectural alignment is a stable property of the data (degree-3), while individual-level detection is measurement-dependent (degree-2).

\textbf{Fatigue invariance.} Dortmund's pre/post design (separated by a 2-hour cognitive battery) tests whether fatigue alters lattice architecture. $\phisym$'s degree-3 $\bar{d}/\bar{d}_{\text{null}}$ changes from 0.914 (EC-pre) to 0.909 (EC-post)---effectively invariant. $\phisym$ maintains rank 1/9 in both EC conditions. Individual-level $d$ drops from 0.391 to 0.265, consistent with the measurement-sensitivity interpretation: post-task spectral peaks are slightly broader (lower SNR), reducing individual-level detectability without altering the underlying architecture.

\subsection{Condition-dependent position migration}
\label{sec:ec_theta}

Eyes-open and eyes-closed resting states produce well-characterized spectral differences \parencite{Barry2007}. We test whether these state changes reorganize lattice position occupancy.

\subsubsection{Architecture is partially stable, partially dynamic}

Using matched extraction methods (both standard [1,\,45]~Hz FOOOF), $N = 182$ subjects with both EO and EC (Table~\ref{tab:position_stability}):

\begin{table}[H]
\centering
\caption{EO--EC position stability by band.}
\label{tab:position_stability}
\begin{tabular}{lccc}
\toprule
Band & Position match (EO $=$ EC) & Freq $r$(EO, EC) & Mean $|\Delta u|$ \\
\midrule
Delta & 31\% & 0.100 & 0.225 \\
Theta & 74\% & 0.237 & 0.153 \\
Alpha & 41\% & 0.388 & 0.174 \\
Gamma & 35\% & $-0.056$ & 0.225 \\
\bottomrule
\end{tabular}
\end{table}

Theta is the most stable band; delta and gamma are essentially uncorrelated across conditions.

\subsubsection{Directional theta convergence on $\fzero$}

Theta position migration is \textbf{asymmetric and directional} (Table~\ref{tab:theta_transition}).

\begin{table}[H]
\centering
\caption{Theta position transition matrix (EO $\to$ EC), using matched extraction ($N = 182$).}
\label{tab:theta_transition}
\begin{tabular}{lcccc}
\toprule
EO position & $N$ & Stay (\%) & $\to$ Boundary (\%) & $\to$ Other (\%) \\
\midrule
Boundary & 128 & 92 & --- & 8 \\
Noble$_1$ & 41 & 32 & \textbf{61} & 7 \\
Noble$_2$ & 6 & 50 & 17 & 33 \\
Attractor & 7 & 14 & 43 & 43 \\
\bottomrule
\end{tabular}
\end{table}

Boundary is an absorbing state: 92\% of subjects with theta at $\fzero$ in EO remain there in EC, while 61\% of noble$_1$ subjects migrate to boundary. The migration is toward $\fzero$ from both directions:
\begin{itemize}
    \item Subjects with EO theta below $\fzero$ ($N = 128$): mean shift $= +0.38$~Hz (upward toward $\fzero$)
    \item Subjects with EO theta above $\fzero$ ($N = 54$): mean shift $= -0.22$~Hz (downward toward $\fzero$)
    \item $|\theta - \fzero|$ decreases from 0.71~Hz (EO) to 0.53~Hz (EC), paired $t = -3.27$, $p = 0.001$
\end{itemize}

Noble$_1 \to$ boundary shifters ($N = 25$) land at EC theta median $= 7.76$~Hz, with 22/25 within 0.5~Hz of $\fzero = 7.83$.

\textbf{Dortmund replication (overlap-trim extraction).} The LEMON analysis above used matched standard extraction (Section~\ref{sec:ec_theta}); Dortmund's replication uses overlap-trim extraction because the standard-extraction theta median in Dortmund falls at 5.00~Hz (see Section~\ref{sec:theta_convergence}), too far below $\fzero$ for standard FOOOF to resolve the convergence cleanly. Under overlap-trim, the same convergence pattern appears ($N = 608$): mean per-subject $|\theta - \fzero|$ decreases from 1.49~Hz (EO-pre) to 0.95~Hz (EC-pre), $d = 0.39$ ($p = 5.6 \times 10^{-20}$). (Note: the 1.49~Hz value is the mean of individual $|\theta_i - \fzero|$ distances, not $|\text{median}(\theta) - \fzero|$; the overlap-trim population median is 6.56~Hz, giving $|\text{median} - \fzero| = 1.27$~Hz.) Boundary occupancy increases from 36\% to 54\% under EC. The migration is directional: below-$\fzero$ subjects shift upward ($+0.755$~Hz), above-$\fzero$ subjects shift downward ($-0.461$~Hz). Position stability is lower in Dortmund (32\% vs LEMON's 74\%), but boundary positions are 63\% stable---consistent with the absorbing-state interpretation.

\textbf{Leakage control.} EC theta-alpha power correlation ($r = 0.448$) is similar to EO ($r = 0.473$), ruling out EC-specific alpha tail leakage as the mechanism.

\subsubsection{Target disambiguation: $\fzero$, not an IAF subharmonic}
\label{sec:disambiguation}

\textbf{LEMON partial-correlation evidence.} IAF varies from 8.66 to 11.70~Hz across LEMON subjects (mean $= 10.09 \pm 0.56$; consistent with the range reported by \cite{Haegens2014}). The partial correlation between EC theta frequency and IAF, controlling for EO theta baseline, is $r = -0.213$ ($p = 0.004$). The \textit{negative} direction rules out all linear IAF-subharmonic hypotheses, which predict positive partials. Furthermore, EC alpha and theta shifts are uncorrelated ($r = 0.072$, $p = 0.33$), ruling out direct alpha-to-theta entrainment.

The fixed $\fzero = 7.83$ target achieves the lowest mean prediction error ($|\theta_{\text{EC}} - \text{target}| = 0.531$~Hz), beating all IAF-derived alternatives (all Wilcoxon $p < 10^{-12}$ vs $\fzero$; Cliff's $\delta = 0.66$, 83\% of subjects closer to $\fzero$). IAF-stratified analysis confirms IAF-invariant convergence ($r = -0.011$, $p = 0.88$; Table~\ref{tab:iaf_tercile}).

\begin{table}[H]
\centering
\caption{IAF-stratified theta convergence (LEMON).}
\label{tab:iaf_tercile}
\begin{tabular}{lccccccc}
\toprule
Tercile & $N$ & Mean IAF & $\theta_{\text{EO}}$ & $\theta_{\text{EC}}$ & $\Delta\theta$ & Conv.\ on $\fzero$ & IAF/$\sqrt{\phisym}$ \\
\midrule
Low IAF & 61 & 9.47 & 7.47 & 7.74 & $+0.27$ & $+0.23^*$ & 7.45 \\
Mid IAF & 60 & 10.07 & 7.24 & 7.38 & $+0.15$ & $+0.14$ & 7.92 \\
High IAF & 61 & 10.71 & 6.83 & 7.02 & $+0.20$ & $+0.18$ & 8.42 \\
\bottomrule
\end{tabular}
\end{table}

However, in LEMON, $\fzero$ and IAF/2 lie in the same direction relative to most theta peaks (both are above the population theta median), so the negative partial rules out linear IAF scaling but cannot distinguish $\fzero$ from a nonlinear ceiling effect.

\textbf{Dortmund sign-flip: definitive disambiguation.} The Dortmund dataset resolves this ambiguity through a geometric accident: with lower IAFs (mean $\sim$10.0~Hz) and overlap-trim EO theta at 6.56~Hz, the competing targets lie in \textit{opposite directions}---$\fzero = 7.83$~Hz is above the theta peak, while IAF/2 $\approx 5.0$~Hz is below it. Under EC, theta shifts \textit{upward} (Figure~\ref{fig:signflip}):

\begin{itemize}
    \item $|\theta - \fzero|$ decreases: $d = -0.385$ ($p = 5.6 \times 10^{-20}$)---theta converges on $\fzero$
    \item $|\theta - \text{IAF}/2|$ \textit{increases}: $d = +0.351$ ($p = 4.3 \times 10^{-17}$)---theta diverges from IAF/2
\end{itemize}

The sign flip is definitive. EC theta moves toward $\fzero$ and away from IAF/2, simultaneously. This excludes all IAF-derived models---linear and nonlinear---because any model where the theta target is a monotonic function of IAF predicts that EC theta should move toward IAF/2 (or at minimum not diverge from it). The $\fzero$ convergence survives alpha-frequency partialling ($t = 9.51$, $p = 4.2 \times 10^{-20}$), and the theta-alpha shift correlation is negligible ($r = -0.082$), confirming independent mechanisms.

\textbf{Robustness.} In LEMON, the convergence-direction regression ($\Delta\theta$ on target gap) yields $R^2 = 0.215$ for $\fzero$ vs 0.129 for IAF/$\sqrt{\phisym}$ (Williams test: $t = 4.24$, $p < 0.0001$). Three IAF estimators all produce negative partial correlations (see Appendix~\ref{app:disambiguation_robustness}). Age stratification confirms $\fzero$ is closer for 83\% of subjects in both young and elderly subgroups.

\begin{keyinsight}
Two datasets, opposite geometric configurations, same conclusion: EC theta converges on the fixed frequency $\fzero = 7.83$~Hz and diverges from IAF/2. The Dortmund sign flip ($d = -0.385$ toward $\fzero$, $d = +0.351$ away from IAF/2) definitively excludes all IAF-derived convergence models, both linear and nonlinear. The convergence target is a fixed frequency, not a function of individual spectral architecture.
\end{keyinsight}

\subsubsection{Power changes at stable positions}

At positions that remain stable across conditions, power changes are large (Table~\ref{tab:power_change}).

\begin{table}[H]
\centering
\caption{Power changes at stable lattice positions (EC vs EO, matched extraction).}
\label{tab:power_change}
\begin{tabular}{lccc}
\toprule
Band & $d$(EC $-$ EO) & $p$ & $N$ (stable) \\
\midrule
Theta & $+0.58$ & $<0.0001$ & 135 \\
Alpha & $+1.00$ & $<0.0001$ & 69 \\
Gamma & $+0.14$ & 0.022 & 96 \\
Delta & $-0.13$ & 0.16 & 60 \\
\bottomrule
\end{tabular}
\end{table}

The alpha power increase at stable positions ($d = 1.00$) is the largest effect size in this paper---a full standard deviation of amplification at lattice positions that do not move. This confirms a dual mechanism: positions can migrate (theta noble $\to$ boundary), and power at stable positions changes independently (alpha amplification in EC). The lattice is neither fully fixed nor fully dynamic---it is a potential landscape where some wells (boundary/$\fzero$) are deeper than others. The depth of the well determines both positional stability and power amplification under increased oscillatory drive.

\subsubsection{EC strengthens four-band lattice alignment}
\label{sec:ec_fourband}

The theta convergence (Section~\ref{sec:disambiguation}) was discovered using standard extraction, which cannot resolve the other three bands cleanly. To test whether EC strengthens alignment \textit{across all bands}, we extracted overlap-trim peaks for the EC condition ($N = 203$; one additional subject had usable EC but not EO data) and computed the dominant-peak metric under all four extraction$\times$condition combinations (Table~\ref{tab:fourway}).

\begin{table}[H]
\centering
\caption{Four-way dominant-peak alignment: extraction method $\times$ condition. Cohen's $d$ reported as unsigned effect size (alignment strength). $^{\dagger}$Anti-aligned ($\bar{d} > \bar{d}_{\text{null}}$, not significant).}
\label{tab:fourway}
\begin{tabular}{lcccc}
\toprule
Metric & OT-EO & OT-EC & STD-EO & STD-EC \\
\midrule
Degree-2 Cohen's $d$ & $0.40$ & $\mathbf{0.75}$ & $0.11^{\dagger}$ & $0.14$ \\
Degree-2 $p$ (one-sided) & $2.3 \times 10^{-8}$ & $\mathbf{1.1 \times 10^{-21}}$ & 0.94 & 0.037 \\
Degree-2 Wilcoxon $p$ & $1.4 \times 10^{-8}$ & $\mathbf{6.2 \times 10^{-18}}$ & 0.83 & 0.010 \\
Mean $\bar{d}$ & 0.069 & \textbf{0.061} & 0.083 & 0.077 \\
\midrule
Degree-3 $\bar{d}/\bar{d}_{\text{null}}$ ($\phisym$) & 0.917 & \textbf{0.874} & 1.020 & 0.902 \\
Degree-3 $\phisym$ rank & 1/9 & \textbf{1/9} & 7/9 & 5/9 \\
Degree-3 $\phisym$ $p$ & $4.8 \times 10^{-5}$ & $\mathbf{8.1 \times 10^{-10}}$ & 0.83 & $1.6 \times 10^{-6}$ \\
Degree-3 $\phisym$ Cohen's $d$ & $0.28$ & $\mathbf{0.45}$ & $0.07^{\dagger}$ & $0.36$ \\
$N$ & 202 & 203 & 196 & 179 \\
\bottomrule
\end{tabular}
\end{table}

\noindent Three results emerge:

\begin{enumerate}
\item \textbf{OT-EC produces the largest effect in this paper.} At $d = 0.75$ ($p = 10^{-21}$), the OT-EC alignment effect is nearly double that of OT-EO ($d = 0.40$). Dominant peaks fall 12.6\% closer to $\phisym$-lattice positions than chance, compared to 8.3\% under OT-EO.

\item \textbf{The interaction is multiplicative.} EC alone provides a modest improvement (STD-EO $\to$ STD-EC: $d$ goes from $0.11$ to $0.14$). Overlap-trim alone provides a large improvement (STD-EO $\to$ OT-EO: $d = 0.11$ to $0.40$). Together, the effects combine to $d = 0.75$---larger than the sum of the individual improvements, consistent with EC increasing the signal (stronger oscillations settle more precisely into lattice positions) while overlap-trim increases measurement accuracy (better aperiodic modeling isolates the genuine peaks).

\item \textbf{All four bands tighten under EC.} The per-band mean lattice distance decreases in every band under OT-EC vs OT-EO: theta $-0.013$, alpha $-0.013$, gamma $-0.006$, delta $-0.002$ (paired $p = 0.0005$; EC better for 61\% of subjects). This extends the theta convergence finding: closing the eyes does not only pull theta toward the boundary---it tightens all four band modes toward their respective lattice positions. Per-position analysis in Dortmund reveals that this tightening is band-specific: theta enrichment at the boundary approximately doubles from EO to EC ($+149\%$ to $+312\%$, $z = 7.0$), while alpha enrichment at attractor and noble$_1$ likewise doubles ($+45\%/+25\%$ EO to $+101\%/+70\%$ EC). Gamma enrichment at noble$_1$ follows the same pattern ($+83\%$ EO to $+173\%$ EC). Each puller is amplified through its own band rather than through a uniform signal increase.
\end{enumerate}

\noindent Under standard extraction, EC alone already achieves significant alignment ($d = 0.36$, $p = 1.6 \times 10^{-6}$ at degree 3), while EO remains non-significant ($d = 0.07$, $p = 0.83$). The condition effect is thus partially independent of the extraction method. However, $\phisym$ ranks 1st in LEMON only under overlap-trim; under standard extraction, bases with more positions (1.8, 1.7) outperform $\phisym$. Accurate peak isolation remains necessary for $\phisym$ to be distinguished from competing bases.

\subsection{Independent confirmation of the Schumann anchor}
\label{sec:anchor_confirmation}

The theta convergence target (7.38~Hz best-fit constant) and the multi-band lattice anchor ($\fzero = 7.83$~Hz, pre-specified) are derived from different analyses. If 7.38~Hz were the ``true'' anchor frequency, the 4-band dominant-peak alignment should improve when recomputed at $\fzero = 7.38$. It does not: $\bar{d} = 0.082$ ($p = 0.29$, $d = 0.07$), far worse than $\fzero = 7.83$ ($\bar{d} = 0.069$, $p < 10^{-7}$, $d = 0.40$).

A fine-grained sweep ($\fzero = 7.0$--$8.5$~Hz in 0.05~Hz steps) confirms the 4-band lattice optimum at $\fzero = 7.88$~Hz---within 0.05~Hz of the Schumann fundamental and far from the theta landing zone. The significant window (Cohen's $d > 0.20$) spans $\fzero \in [7.50,\, 8.20]$~Hz, centered on 7.83. Below 7.45~Hz, alignment is non-significant; above 8.35~Hz, it returns to null.

This result is independent of the theta disambiguation: it uses all four bands (not just theta), the overlap-trim extraction (not the standard extraction used for EC/EO comparison), and the pre-specified anchor (no optimization). The 7.38 best-fit constant reflects where theta-band convergence lands (one band); 7.88 is where the multi-band lattice is anchored (four bands). Theta converges \textit{toward} but undershoots the lattice anchor, consistent with the basin interpretation.

A second, independent line of evidence corroborates this offset. The per-band KDE density at the boundary position ($u = 0$) reveals that theta's enrichment peak is not centered at $u = 0$ but at $u \approx +0.02$, corresponding to an effective $\fzero$ slightly above 7.83~Hz ($f_{\text{eff}} = 7.90$--$7.94$~Hz across the three datasets; cross-dataset range 0.04~Hz, below the spectral resolution of any individual dataset). The 4-band lattice sweep (above) and the single-band theta KDE use different data, different statistics, and different definitions of ``best anchor,'' yet converge on an effective frequency 0.05--0.10~Hz above the canonical Schumann fundamental. Whether this offset reflects genuine geophysical variation (the Schumann fundamental varies diurnally by $\pm$0.5~Hz; \cite{Nickolaenko2014}) or a systematic theta-band bias cannot be determined from scalp EEG alone.

\subsection{Cross-dataset EO theta convergence}
\label{sec:theta_convergence}

Under standard extraction, LEMON and Dortmund theta populations start at very different frequencies: LEMON EO median $= 7.64$~Hz (above $\fzero$), Dortmund EO median $= 5.00$~Hz (below $\fzero$). Under overlap-trim extraction, both converge to nearly identical values: LEMON OT-EO $= 6.5685$~Hz, Dortmund OT-EO $= 6.5594$~Hz---a cross-dataset difference of $\Delta = 0.009$~Hz (Figure~\ref{fig:theta_convergence}).

Both values fall at the noble$_1$ position ($u \approx 0.633$, within 0.015 of the theoretical $1/\phisym = 0.618$). The convergence arrives from opposite directions: LEMON theta shifts downward by 1.07~Hz (from 7.64 to 6.57), Dortmund theta shifts upward by 1.56~Hz (from 5.00 to 6.56). Two independent datasets, recorded with different equipment at different laboratories, converge to the same lattice coordinate under the same extraction procedure. The $0.009$~Hz difference is orders of magnitude smaller than the standard deviation of individual theta frequencies ($\sim$1.5~Hz) and well below the spectral resolution of either dataset.

\subsection{Population stability despite individual reshuffling}
\label{sec:pop_stability}

Dortmund's pre/post design reveals that the lattice describes population attractors rather than individual fingerprints. At the population level, lattice architecture is effectively invariant: $\phisym$'s degree-3 $\bar{d}/\bar{d}_{\text{null}}$ shifts from 0.914 (EC-pre) to 0.909 (EC-post)---a negligible change. $\phisym$ remains rank 1/9 in both EC conditions.

At the individual level, however, pre-post alignment is weak: per-subject $\bar{d}$ correlates $r = 0.131$ between pre and post sessions. Individuals reshuffle across lattice positions while the population-level occupancy remains stable. This is analogous to a crystal lattice: the structure describes allowed sites; which sites are occupied by specific atoms fluctuates. The five-year longitudinal arm confirms this is genuine instability, not measurement noise: ICC(2,1) is negative across all conditions (Section~\ref{sec:longitudinal}). The 15 position-pair combinations documented in Section~\ref{sec:phenotyping} reflect this individual diversity within a stable population scaffold.

\subsection{Five-year longitudinal stability: population architecture, not individual trait}
\label{sec:longitudinal}

Section~\ref{sec:pop_stability} introduced the crystal-lattice analogy: sites are stable, occupants fluctuate. But that evidence spanned only a 2-hour interval, leaving open whether the individual reshuffling ($r = 0.131$) reflects genuine instability or short-term measurement noise. The Dortmund longitudinal arm---208 subjects returning approximately five years after session~1 for an identical protocol (four resting-state blocks: EC-pre, EO-pre, EC-post, EO-post)---upgrades this analogy from minutes to years and provides a direct test of whether $\phisym$-alignment is a stable individual trait or an emergent property of population-level frequency architecture. To ensure matched within-subject comparisons, this section uses standard [1,\,45]~Hz FOOOF extraction for both sessions (the overlap-trim method was applied only to ses-1 in the main analysis); absolute $\bar{d}/\bar{d}_{\text{null}}$ values therefore differ slightly from the overlap-trim results reported elsewhere.

\subsubsection{Population-level invariance across five years}

At the population level, lattice architecture is effectively a constant of the system. Group mean $\bar{d}$ is locked at 0.031--0.033 across all eight condition$\times$session cells (Table~\ref{tab:longitudinal}). In ses-2, $\phisym$ ranks 1st in one condition (EO-post, $\bar{d}/\bar{d}_{\text{null}} = 0.934$), 2nd in EO-pre (behind $e$), 3rd in EC-pre (behind 1.8 and $e$), and 6th in EC-post ($\bar{d}/\bar{d}_{\text{null}} = 1.009$, not significant). This variability, compared with the stable rank 1--2 in ses-1, likely reflects the smaller sample ($N = 208$ vs 608) and the narrower gap between competing bases. Under eyes-open conditions, ses-2 alignment reaches standalone significance: $d = 0.245$, $p < 0.001$ (pre) and $d = 0.230$, $p = 0.001$ (post)---effect sizes consistent with the full $N = 608$ ses-1 analysis.

\begin{table}[H]
\centering
\caption{Longitudinal stability: ses-1 vs ses-2 ($N = 208$ matched subjects).}
\label{tab:longitudinal}
\begin{tabular}{lccccccc}
\toprule
Condition & $\bar{d}_{\text{s1}}$ & $\bar{d}_{\text{s2}}$ & $\Delta$ & Paired $p$ & Pearson $r$ & ICC(2,1) & Ses-2 $d$ \\
\midrule
EC-pre  & 0.0322 & 0.0325 & $+$0.0002 & 0.787 & 0.008 & $-$0.33 & $-$0.02 \\
EO-pre  & 0.0315 & 0.0303 & $-$0.0013 & 0.129 & 0.095 & $-$0.25 & 0.25 \\
EC-post & 0.0319 & 0.0324 & $+$0.0005 & 0.543 & $-$0.009 & $-$0.34 & $-$0.01 \\
EO-post & 0.0314 & 0.0303 & $-$0.0010 & 0.250 & $-$0.034 & $-$0.36 & 0.23 \\
\bottomrule
\end{tabular}
\end{table}

\subsubsection{Individual-level alignment is not a trait}

Within-subject stability across five years is null. Pearson correlations between ses-1 and ses-2 $\bar{d}$ range from $r = -0.034$ to $r = 0.095$; all ICC(2,1) values are negative ($-0.25$ to $-0.36$). Negative ICCs indicate that between-subject variance is \textit{smaller} than within-subject variance across timepoints---the defining signature of a non-trait measure. Mean absolute change $|\Delta\bar{d}| \approx 0.010$ is small in absolute terms (roughly one-third of the population mean $\bar{d}$), but individual alignment at ses-1 has no predictive value for alignment at ses-2.

This resolves the ambiguity flagged in Section~\ref{sec:pop_stability}: the within-session pre-post correlation ($r = 0.131$) could have reflected either genuine instability or measurement noise over a short interval. The five-year ICC confirms that the instability is real---individuals reshuffle across lattice positions on timescales from minutes to years, while the population-level distribution remains invariant.

\subsubsection{Per-band position stability gradient}

Per-band position stability tracks the physiological stability of the underlying oscillator. Alpha is most stable (33--47\% same nearest position across sessions, frequency $r = 0.46$--$0.66$), consistent with individual alpha frequency being one of the most heritable EEG features. Theta shows moderate stability (23--43\%, frequency $r = 0.35$--$0.43$). Delta and gamma are unstable (15--27\%, frequency $r = 0.08$--$0.34$). The lattice positions inherit whatever test-retest reliability the underlying oscillatory generators possess, but the lattice \textit{relationship}---the population-level convergence on $\phisym$-positions---is a constant.

\subsubsection{Age does not moderate stability}

Alignment stability is age-invariant: $|\Delta\bar{d}|$ correlates $r = -0.03$ to $-0.08$ with age across conditions (all $p > 0.27$). Young ($<$35) and elderly ($\geq$55) subjects show indistinguishable absolute change ($|\Delta\bar{d}| \approx 0.010$ in both groups). The architectural property neither strengthens nor degrades with aging.

\subsubsection{Fatigue invariance replicates in session 2}

Pre-post cognitive battery differences in ses-2 replicate the ses-1 null: $\Delta\bar{d} < 0.001$ in all four condition$\times$session cells. This confirms that alignment is insensitive to acute cognitive fatigue at both timepoints, ruling out state confounds in the longitudinal comparison.

\subsubsection{Interpretation}

The dissociation between population-level invariance (group $\bar{d}$ locked, $\phisym$ rank 1--3 across conditions) and individual-level instability (ICC $< 0$) identifies lattice alignment as a species-level architectural constant rather than an individual biomarker. The crystal-lattice analogy from Section~\ref{sec:pop_stability} extends from minutes to years: the allowed sites are stable; which sites a given brain occupies at any moment fluctuates. This is consistent with the cognitive null (Section~\ref{sec:cognitive})---a population-level structural constant would not be expected to predict individual cognitive performance, any more than a structural constant of the medium would be expected to predict variation in signals transmitted through it. However, this reframing is a post-hoc interpretation; the ICC result was not predicted by the cognitive null, and the cognitive null (tested only in LEMON) would require direct replication in a dataset with both longitudinal and cognitive data to confirm the link.

\subsection{Extraction sensitivity}
\label{sec:phenotyping}

Under standard whole-range FOOOF [1,\,45]~Hz, $\bar{d} = 0.083$ ($d = 0.12$, $p = 0.08$); under overlap-trim, $\bar{d} = 0.069$ ($d = 0.40$, $p < 10^{-7}$). One-third of subjects change theta position between methods, confirming that position assignment is partly extraction-dependent. The within-Dortmund condition gradient (Section~\ref{sec:measurement}) shows this is measurement precision, not architecture, varying across methods and conditions. All dominant-peak results use overlap-trim extraction; all EO--EC comparisons use matched methods. Position-type phenotyping (58\% boundary, 25\% noble$_1$; 15 theta-alpha position-pair combinations; no cognitive differences between groups) appears in Table~S4.

% =====================================================================
\section{Discussion}
% =====================================================================

\subsection{The diagnostic result: aggregate metrics are methodological}

Five diagnostic approaches converge on a single conclusion: all-peaks enrichment metrics do not measure neural architecture. They measure the interaction between FOOOF extraction parameters and the lattice coordinate system.

The cascade:
\begin{enumerate}
    \item Standard SS partially replicates but depends on frequency range (collapses at [1,\,85]~Hz)
    \item Amplitude weighting attenuates rather than strengthening age correlations
    \item Per-$\phisym$-octave fitting produces enrichment that tracks fitting window edges, not lattice positions
    \item The age effect decomposes into alpha+beta IAF positioning; per-subject $\fzero$ correction reverses its sign
    \item Without an anchor frequency, $\phisym$ ranks 4th among nine bases
\end{enumerate}

Each step eliminated a potential rescue for the previous finding. This progression is the methodological contribution of the paper: it demonstrates that spectral peak enrichment analysis requires controls for extraction parameters \parencite{Donoghue2022methods, Ostlund2022}, anchor sensitivity, and anchor-free alternatives that the current literature has not applied.

\subsection{The constructive result: dominant-peak alignment replicates across three datasets}

The dominant-peak metric is substantially more robust than all-peaks metrics: one peak per band, conventional band boundaries, pre-specified anchor. The effect ($d = 0.40$, $p < 10^{-7}$ under overlap-trim extraction at degree 2) is present in individual brains, age-invariant, fatigue-invariant, and replicates across three independent datasets totaling 919 subjects ($\bar{d} = 0.069$ (range: 0.0689--0.0693)), including a five-year longitudinal retest of 208 subjects. Under coverage normalization with degree-3 noble positions, $\phisym$ ranks 1st in all three datasets: LEMON ($\bar{d}/\bar{d}_{\text{null}} = 0.917$, $p < 10^{-4}$), Dortmund ($0.914$, $p < 10^{-11}$), and EEGMMIDB ($0.901$, $p < 10^{-3}$), achieving 8--10\% closer alignment than chance (Table~\ref{tab:crossbase}). The second-place base varies across datasets (1.7, 3/2, $\sqrt{2}$), indicating that $\phisym$'s rank-1 consistency is not shared by any competitor.

The metric is not fully extraction-immune ($d = 0.12$, $p = 0.08$ under standard extraction), but the dependence is qualitatively different from the all-peaks case. All-peaks metrics \textit{change sign} with parameter variation; the dominant-peak metric changes \textit{magnitude} but not direction. Moreover, per-octave aperiodic fitting is the more principled method: FOOOF's single-exponent model poorly approximates the full 1--45~Hz aperiodic spectrum, which exhibits a spectral knee near 10--20~Hz \parencite{Donoghue2020, Gao2017}. Per-octave fitting was adopted as the methodological fix for Part~I's diagnostic finding, not because it produced a favorable $\phisym$ result. The $\phisym$-specificity claim is thus conditional on per-octave aperiodic modeling; under standard whole-range extraction, $\phisym$ does not separate from competing bases (rank 7/9, $p = 0.83$). This conditionality is the single most important caveat in the paper, and it has a clear resolution: independent replication using IRASA \parencite{Wen2016} or eBOSC \parencite{Kosciessa2020}, which separate periodic from aperiodic components through fundamentally different algorithms (resampling-based and wavelet-based, respectively, vs.\ FOOOF's parametric fit). If $\phisym$ maintains its rank-1 position under IRASA-extracted dominant peaks with no per-octave fitting, the alignment signal becomes method-independent. If it does not, the lattice alignment is an artifact of the aperiodic correction step. We identify this as the highest-priority replication target.

The position-count decomposition (Section~\ref{sec:crossbase}) provides independent validation: $d$ attenuates monotonically from $\sim$0.40 (4 positions) to $\sim$0.28 (6 positions) to $\sim$0.07 (8 positions) across all three datasets, as expected for genuine signal diluted by noise positions. This creates a sharp separation: the lattice describes \textit{where} the four strongest oscillatory generators sit in frequency space (a genuine architectural property), but \textit{not} how the broader spectral landscape organizes (an extraction artifact) or how effectively those generators drive cognition (a null). Only three of six degree-3 positions carry signal in the aggregate cross-band analysis (noble$_1$, attractor, boundary); individual frequency bands may exhibit distinct position preferences, consistent with the broader finding that aggregate all-peaks metrics conflate band-specific patterns.

\subsection{The EC theta migration: a potential landscape confirmed by disambiguation}

The directional convergence of theta toward $\fzero$ under eyes-closed conditions reveals the lattice as a potential landscape rather than a rigid grid. The boundary position at $\fzero$ is the deepest well: 92\% stable in LEMON, absorbing 61\% of noble$_1$ occupants when oscillatory drive increases. Dortmund replicates the pattern ($d = -0.39$, $N = 608$). Noble positions depopulate under increased drive, suggesting they represent metastable states that require active maintenance.

The Dortmund sign-flip definitively resolves the convergence target (Section~\ref{sec:disambiguation}). In LEMON, $\fzero$ and IAF/2 both lie above the theta peak, so partial correlations could distinguish linear but not nonlinear IAF-derived models. In Dortmund, $\fzero$ and IAF/2 lie in \textit{opposite} directions from theta---and EC theta moves toward $\fzero$ ($d = -0.385$) while simultaneously diverging from IAF/2 ($d = +0.351$). This sign flip excludes all IAF-derived convergence models, linear and nonlinear. The convergence target is a fixed frequency at $\fzero = 7.83$~Hz. The cross-dataset theta medians converge from opposite directions (LEMON: 7.64 $\to$ 6.57~Hz; Dortmund: 5.00 $\to$ 6.56~Hz) to land 0.009~Hz apart under overlap-trim extraction---a striking coincidence that further supports a shared architectural attractor.

The dual mechanism (position migration + power amplification at stable positions) bridges the fixed-architecture and dynamic-redistribution framings: the lattice provides the positions, conditions select the occupancy, and drive modulates the amplitude.

The four-band EC analysis (Section~\ref{sec:ec_fourband}) extends this interpretation beyond theta. Under overlap-trim extraction, EC strengthens alignment across \textit{all four bands} ($d = 0.75$ vs $0.40$ for EO; $p = 10^{-21}$), with theta and alpha contributing equally ($\Delta d = 0.013$ each). The interaction between extraction method and condition is multiplicative, consistent with EC increasing the underlying signal while overlap-trim increases measurement accuracy.

The per-position band decomposition clarifies the mechanism. Each of the three signal-carrying positions is driven by a different band: theta dominates the boundary ($+312\%$ enrichment, $5{\times}$ stronger than any other band at that position), alpha and gamma jointly drive attractor and noble$_1$ (alpha $+101\%/+70\%$; gamma $+173\%$ at noble$_1$), and delta contributes no consistent positional signal. EC amplifies each position through its respective band: theta doubles at boundary, alpha doubles at attractor/noble$_1$, gamma doubles at noble$_1$---a band-specific multiplicative effect, not a generic signal increase. Crucially, the theta boundary offset documented in Section~\ref{sec:anchor_confirmation} (KDE peak at $u = +0.02$, effective $\fzero \approx 7.92$~Hz) is condition-invariant: EC doubles the amplitude of the boundary cluster without shifting its location. The lattice positions are architectural constants; the state modulation is purely energetic.

This decomposition makes a testable prediction: positions in spectral gaps between band-specific attractors should be depleted regardless of condition, and the depletion pattern should shift only when the band landscape changes. The Dortmund EO/EC comparison confirms this. Inv-noble$_4$ ($u = 0.854$) is the only position that changes sign between conditions ($-16\%$ under EC, $+18\%$ under EO), traceable to theta redistribution: in EO, theta's weaker boundary concentration spreads density into the $u = 0.85$ region ($+34\%$ theta enrichment in EO vs $-6\%$ in EC). This exception-that-proves-the-rule specificity distinguishes the band-decomposition model from a generic ``EC amplifies everything'' account.

One novel observation: gamma enrichment at noble$_1$ ($u = 0.618$, the $1/\varphi$ position) doubles from $+83\%$ in EO to $+173\%$ in EC. Gamma-band organization at $\varphi$-lattice positions being state-modulated is not predicted by existing EO/EC literature, and deserves targeted investigation with broader gamma bandwidth ($>45$~Hz) to determine whether this reflects genuine gamma-$\varphi$ coupling or a bandwidth ceiling artifact (Section~\ref{sec:limitations}).

\subsection{Cognitive silence across three metrics within LEMON}

\begin{table}[H]
\centering
\caption{Cognitive prediction null across all metrics.}
\begin{tabular}{lccc}
\toprule
Metric & Tests & FDR survivors & Note \\
\midrule
All-peaks enrichment (SS) & 9 & 0 & --- \\
Within-band $z$-scores & 36 & 1 & Non-$\phisym$-specific \\
Dominant-peak distance ($\bar{d}$) & 36+ & 0 & --- \\
\bottomrule
\end{tabular}
\end{table}

The lattice does not predict cognitive performance---and this single-dataset null across three independent metrics is itself a substantive finding that demands explanation. Three metrics applied to LEMON data across 100+ tests yield zero FDR survivors. If the $\phisym$-lattice reflects genuine neural architecture, as three-dataset convergence suggests, then the organizational principle governing frequency placement is decoupled from the computational function of those oscillations \parencite{Siegel2012}. Dortmund's fatigue-invariance sharpens the challenge: a 2-hour cognitive battery leaves degree-3 architecture unchanged ($\bar{d}/\bar{d}_{\text{null}}$: 0.914 $\to$ 0.909), ruling out state-dependent cognitive sensitivity alongside trait-dependent prediction.

Three caveats temper this conclusion. First, cognitive prediction is tested only in LEMON ($N = 202$), which provides 80\% power only down to $|r| = 0.196$. Small but theoretically meaningful effects ($|r| \sim 0.10$--$0.15$) are not ruled out, and the null is a single-dataset result: neither Dortmund nor EEGMMIDB provides cognitive data. A combined sample of $N = 919$ would detect effects as small as $|r| = 0.09$ at 80\% power---if cognitive measures were available across all three datasets, the null would be substantially more constraining. Second, Dortmund's fatigue-invariance ($\bar{d}/\bar{d}_{\text{null}}$: 0.914 $\to$ 0.909 across a 2-hour cognitive battery) provides indirect support for cognitive independence, but it tests state sensitivity, not trait-level prediction---it is not a substitute for direct cognitive testing. Third, the longitudinal result (Section~\ref{sec:longitudinal}) is consistent with the cognitive null: if $\phisym$-alignment is a population-level constant with no individual test-retest reliability (ICC $< 0$), then individual cognitive correlations are structurally unlikely---a measure that does not reliably distinguish individuals cannot predict individual differences in cognition. However, this is a post-hoc interpretation, not a pre-registered test, and should be treated as a framing rather than evidence.

This dissociation constrains any theoretical account of frequency architecture: explanations that derive cognitive function from oscillatory spacing must account for why a replicable, $\phisym$-specific positioning effect carries no detectable cognitive consequence and shows no individual stability. The scaffolding is real; its functional significance, if any, lies elsewhere.

\subsection{What the noble-position advantage means}

The cross-base comparison (Section~\ref{sec:crossbase}) reveals a subtlety unique to $\phisym$. At degree-2 nobles---the standard choice---$\phisym$'s self-similarity collapses six candidate positions to four, and $\phisym$ ranks last among nine bases. This initially appears to contradict the $\phisym$-lattice hypothesis. In fact, it demonstrates it: the collapse is a direct consequence of $\phisym = 1 + 1/\phisym$, the defining recurrence of the golden ratio.

At degree 3, $\phisym$ gains the positions $u = 0.236$ ($(1/\phisym)^3$) and $u = 0.764$ ($1 - (1/\phisym)^3$), yielding six distinct positions. These are Fibonacci fractions---consecutive terms of the sequence $1/\phisym, (1/\phisym)^2, (1/\phisym)^3, \ldots$---which the three-distance theorem guarantees to be optimally equidistributed on the unit interval \parencite{Slater1967}. Under coverage-normalized comparison ($\bar{d}/\bar{d}_{\text{null}}$), $\phisym$ ranks 1st in all three datasets ($0.901$--$0.917$), consistent with $\phisym$'s optimal equidistribution conferring a genuine advantage. However, the margins are narrow in LEMON (0.917 vs 1.7's 0.919) and Dortmund (0.914 vs 3/2's 0.919), and the ranking is extraction-dependent (Section~\ref{sec:crossbase}).

The per-position analysis (Table~\ref{tab:per_position}) reveals a subtlety that tempers this result: the degree-3 positions themselves carry no signal. Peaks are significantly \textit{anti-aligned} with $u = 0.236$ and $u = 0.764$ across all three datasets (Table~\ref{tab:deg3_critical}). $\phisym$'s rank-1 advantage at degree 3 therefore derives from the coverage-normalization properties of its position \textit{spacing}---its six positions tile the unit interval most uniformly, yielding the lowest $\bar{d}_{\text{null}}$ among six-position bases---not from signal at the degree-3 positions themselves. This distinction matters: $\phisym$ is the best lattice because its mathematical structure creates the fairest comparison metric, and under that fair metric, the three aggregate signal-carrying positions (boundary, attractor, noble$_1$) happen to serve $\phisym$ best. But the degree-3 positions are geometric scaffolding for the comparison, not sites of biological alignment.

Notably, only three of the four degree-2 positions carry signal. Noble$_2$ ($u = 0.382$) is consistently null or depleted despite being a core lattice position. The effective signal-carrying subset $\{0, 0.5, 0.618\}$ has a natural geometric interpretation: boundary (octave edge), attractor (geometric midpoint), and the upper noble ($1/\phisym$). Whether noble$_2$'s absence reflects competitive suppression from the nearby attractor (spacing $= 0.118$, vs noble$_1$'s 0.118 from attractor in the other direction) or a genuine three-position architecture requires further investigation. This asymmetry does not affect the cross-base comparison (which uses all six degree-3 positions for all bases) but constrains any theoretical account of why $\phisym$-positions are preferred: the mechanism must explain signal at three positions, not four.

This result is methodologically stronger than the original degree-2 comparison in two ways: (1)~it uses a fair position count (6 positions for $\phisym$, 6--8 for others), and (2)~it controls for position spacing through coverage normalization. Among bases with exactly six positions ($\phisym$, $\sqrt{2}$, 2, 1.4), $\phisym$ significantly outperforms all others ($p < 0.001$), ruling out a generic coverage explanation among equal-count bases. $\phisym$ is the only base that consistently ranks 1st across all three datasets; the second-place base varies (1.7, 3/2, $\sqrt{2}$), suggesting that $\phisym$'s optimal equidistribution provides a unique structural advantage, though with narrow margins.

However, this advantage is detectable only in dominant-peak alignment. In aggregate metrics, it is drowned by extraction artifacts. In pairwise ratios (no anchor), $\phisym$ ranks 4th. The $\phisym$-specific signal is narrow: it describes the four strongest oscillatory generators \parencite{Buzsaki2004}, at a specific anchor, using positions unique to $\phisym$ \parencite{Pletzer2010, Kramer2022, Izhikevich1997}. Extending the claim beyond these bounds produces artifacts.

\subsection{Cross-dataset replication}

\begin{table}[H]
\centering
\caption{Three-dataset comparison. Dortmund values use EC-pre as the primary condition.}
\begin{tabular}{lccc}
\toprule
Finding & EEGMMIDB ($N = 109$) & LEMON ($N = 202$) & Dortmund ($N = 608$) \\
\midrule
Per-subject $\bar{d}$ (4-pos) & 0.0693 & 0.0689 & 0.0690 \\
$\bar{d}/\bar{d}_{\text{null}}$ (deg-3) & 0.901 (rank 1/9) & 0.917 (rank 1/9) & 0.914 (rank 1/9) \\
Cohen's $d$ (deg-2) & 0.377 & 0.400 & 0.391 \\
Age-invariance ($r$) & --- & $+0.054$ (gap) & $-0.004$ (continuous) \\
f$_0$ disambiguation & Ambiguous & Partial (neg partial) & Definitive (sign flip) \\
EC theta $\to$ $\fzero$ & Not tested & $d = -0.39$ (same dir.) & $d = -0.39$ (opp.\ dir.) \\
OT EO theta median & --- & 6.569~Hz & 6.559~Hz \\
Fatigue-invariance & --- & --- & 0.914 $\to$ 0.909 \\
Individual pre-post $r$ & --- & --- & 0.131 \\
5-year ICC(2,1) & --- & --- & $-0.25$ to $-0.36$ \\
5-year $\phisym$ rank (ses-2) & --- & --- & 1st--6th (varies by condition) \\
5-year ses-2 $d$ (EO) & --- & --- & 0.245 ($p < 0.001$) \\
Cognitive prediction & Not tested & Null (0/36+ FDR) & Not available \\
\bottomrule
\end{tabular}
\end{table}

The dominant-peak alignment test replicates across all three datasets ($N = 919$ total; \cite{Babayan2019, Schalk2004, Goldberger2000, GudiMindermann2024}): under coverage normalization with degree-3 symmetric noble positions, $\phisym$ ranks 1st in all three datasets ($\bar{d}/\bar{d}_{\text{null}} = 0.901$--$0.917$, all $p < 10^{-3}$), with peaks falling 8--10\% closer to lattice positions than chance. The per-subject mean lattice distance converges to $\bar{d} = 0.069$ (range: 0.0689--0.0693) across three laboratories, equipment, and preprocessing pipelines. We note that EEGMMIDB served as the analysis dataset for our prior work \parencite{Lacy2026paper2}; only Dortmund is a fully out-of-sample replication. Dortmund's sign-flip definitively resolves the $\fzero$ vs IAF/2 ambiguity that neither LEMON nor EEGMMIDB could cleanly separate. The cross-dataset EO theta convergence ($\Delta = 0.009$~Hz) provides striking evidence for a shared architectural attractor at the noble$_1$ position. The five-year longitudinal arm (Section~\ref{sec:longitudinal}) extends replication across time, with group $\bar{d}$ invariant to the third decimal place, though $\phisym$'s cross-base ranking is more variable in the smaller ses-2 sample (rank 1--6 across conditions).

\subsection{Implications for future $\phisym$-lattice research}

\begin{enumerate}
    \item \textbf{Aggregate metrics require extraction controls.} Any study reporting enrichment at lattice positions must demonstrate stability across FOOOF frequency range, \texttt{max\_n\_peaks}, and fitting window placement. We recommend the overlap-trim method as a minimum standard, with offset controls \parencite{Donoghue2020, Ostlund2022}.
    \item \textbf{Anchor frequency must be pre-specified or eliminated.} Post-hoc $\fzero$ optimization conflates IAF positioning with lattice structure. Either pre-specify $\fzero$ (as we do at 7.83~Hz) or use anchor-free methods (pairwise ratios).
    \item \textbf{Dominant-peak alignment is the viable metric.} It shows significant alignment for $\phisym$ and nearby bases through degree-3 noble positions under coverage normalization ($\phisym$ rank 1/9 across three datasets, $N = 919$). But it answers a narrow question (4 peaks per subject, reflecting canonical network generators; \cite{Buzsaki2004}) and cannot be extended to richer descriptions of spectral organization without re-encountering the sensitivity problems. The alignment signal is concentrated at three positions in the aggregate cross-band analysis (boundary, attractor, noble$_1$; per-band patterns may differ); the degree-3 positions added for fair cross-base comparison carry anti-signal (Table~\ref{tab:per_position}). Future cross-base comparisons must use degree $\geq 3$ symmetric nobles (both $(1/b)^k$ and $1 - (1/b)^k$) and coverage normalization to avoid the position-count asymmetry caused by $\phisym$'s self-similarity collapse at degree 2, but should report per-position alignment to distinguish coverage-normalization effects from genuine signal at higher-degree positions.
    \item \textbf{The cognitive null, though tested in only one dataset, constrains the field.} Three independent metrics applied to LEMON ($N = 202$, 80\% power to $|r| = 0.20$) yield zero FDR survivors across 100+ tests. This is a single-dataset result---neither Dortmund nor EEGMMIDB provides cognitive data---but the consistency across three metrics within LEMON makes a false negative unlikely for moderate-or-larger effects. Future work claiming cognitive correlates of lattice architecture must specify which metric, at which $\fzero$, with which extraction parameters, and demonstrate robustness to the alternatives. Testing cognitive prediction in a large dataset ($N > 500$) with both lattice alignment and comprehensive cognitive measures remains an important gap.
\end{enumerate}

\subsection{Limitations}
\label{sec:limitations}

\begin{enumerate}
    \item \textbf{Individual instability is confirmed, not resolved:} The five-year longitudinal arm ($N = 208$, Section~\ref{sec:longitudinal}) confirms that within-subject alignment is non-trait (ICC $< 0$), ruling out measurement noise as an explanation for the within-session pre-post instability ($r = 0.131$). The instability is genuine: individuals reshuffle across lattice positions on timescales from minutes to years. Whether this reflects stochastic fluctuation of oscillatory generators around fixed attractors or genuine architectural drift cannot be determined without denser longitudinal sampling (e.g., weekly recordings).
    \item \textbf{Gamma ceiling:} FOOOF at [1,\,45]~Hz truncates the gamma $\phisym$-octave. The dominant gamma peak (median 39.85~Hz) is near the range boundary, possibly distorted. Intracranial data show gamma extending well above 45~Hz \parencite{Hermes2015}. The per-band decomposition reveals that gamma is the third strongest contributor to noble$_1$ enrichment, and its EC modulation ($+83\% \to +173\%$) suggests state-dependent gamma-lattice coupling. Both findings are constrained by the 45~Hz ceiling: the true gamma contribution to lattice alignment may be underestimated or distorted.
    \item \textbf{Extraction sensitivity:} 34\% of subjects change theta position between standard and overlap-trim FOOOF on the same data. Lattice position assignment is partly method-dependent (Section~\ref{sec:phenotyping}).
    \item \textbf{Dominant-peak metric narrowness:} Four peaks per subject is a severe compression of spectral information. The metric survives precisely because it discards everything sensitive to extraction parameters---but this also limits what it can claim.
    \item \textbf{Theta undershoot:} The best-fit constant for theta convergence is 7.38~Hz rather than 7.83~Hz. The basin interpretation---that EC theta converges \textit{toward} $\fzero$ without reaching it, landing instead at the noble$_1$ position (6.56~Hz) partway up the potential well---is consistent with the data but cannot be distinguished from theta band-edge truncation or genuine asymmetry in the landscape without higher-resolution frequency estimation.
    \item \textbf{Cognitive data limited to one dataset:} Cognitive prediction is tested only in LEMON ($N = 202$, 80\% power to $|r| = 0.20$); neither Dortmund nor EEGMMIDB provides cognitive measures. The consistency across three metrics within LEMON makes a false negative unlikely for moderate-or-larger effects, but small effects ($|r| \sim 0.10$--$0.15$) cannot be excluded. A combined $N = 919$ sample would detect $|r| = 0.09$ at 80\% power if cognitive data were available across all three datasets. Dortmund's fatigue-invariance ($\bar{d}/\bar{d}_{\text{null}}$: 0.914 $\to$ 0.909 across a 2-hour cognitive battery) provides indirect evidence that the alignment does not track cognitive load, but this tests state sensitivity, not trait prediction.
\end{enumerate}

\subsection{Conclusion}

The $\phisym$-lattice hypothesis, when applied to aggregate spectral peak distributions, does not survive progressive diagnostic testing. Every all-peaks metric we examined is dominated by FOOOF extraction parameters, IAF projection, or both.

However, a more fundamental test survives: the single strongest oscillation per canonical frequency band sits closer to logarithmic lattice positions than expected by chance, and this replicates across three datasets totaling 919 subjects (EEGMMIDB from our prior work; LEMON as primary analysis; Dortmund as out-of-sample replication). Per-subject mean lattice distance converges to $\bar{d} = 0.069$ (range: 0.0689--0.0693) across LEMON ($N = 202$), Dortmund ($N = 608$), and EEGMMIDB ($N = 109$). Under coverage normalization with degree-3 symmetric noble positions, $\phisym$ ranks 1st in all three datasets ($\bar{d}/\bar{d}_{\text{null}} = 0.901$--$0.917$, all $p < 10^{-3}$). The alignment is age-invariant ($r = -0.004$ in Dortmund's continuous 20--70 age range), fatigue-invariant ($0.914 \to 0.909$ across 2-hour cognitive battery), and cognitively silent in the one dataset with available cognitive measures (LEMON: 0/100+ tests survive FDR across three independent metrics).

The Dortmund sign-flip definitively resolves the convergence target: under EC, theta moves toward $\fzero = 7.83$~Hz ($d = -0.385$) and \textit{away from} IAF/2 ($d = +0.351$), excluding all IAF-derived models. Two datasets' EO theta populations converge from opposite starting frequencies (7.64 and 5.00~Hz) to land 0.009~Hz apart at the noble$_1$ lattice position. The position-count decomposition ($d \approx 0.40 \to 0.28 \to 0.07$ from 4 to 6 to 8 positions) is consistent with signal concentrated at the four core lattice positions, though the monotonic attenuation alone cannot distinguish genuine signal dilution from the statistical tendency of coarser partitions to capture larger departures from uniformity (Section~\ref{sec:crossbase}).

The golden ratio describes where dominant oscillatory generators live in frequency space \parencite{Buzsaki2004, Penttonen2003, Pletzer2010}---a genuine architectural property that behaves like a crystal lattice: the allowed sites are stable across datasets, conditions, equipment, and a five-year longitudinal interval, while which sites individual subjects occupy fluctuates (within-session $r = 0.131$; five-year ICC $< 0$). This dissociation identifies lattice alignment as a species-level constant of cortical frequency architecture rather than an individual trait or biomarker. The lattice carries no cognitive consequence we can detect in LEMON ($N = 202$; the only dataset with cognitive measures), but its positions are dynamically meaningful: they are attractors that absorb oscillatory generators under increased drive. The architecture is real, narrow, and inert.

% =====================================================================
\section{Key Statistics Summary}
% =====================================================================

\begin{longtable}{p{4cm}p{3.5cm}p{2.5cm}cl}
\caption{Key statistics summary.} \\
\toprule
Claim & Test & Statistic & $p$ & Survives? \\
\midrule
\endfirsthead
\multicolumn{5}{c}{\textit{Table continued}} \\
\toprule
Claim & Test & Statistic & $p$ & Survives? \\
\midrule
\endhead
\midrule
\multicolumn{5}{r}{\textit{Continued on next page}} \\
\endfoot
\bottomrule
\endlastfoot

\multicolumn{5}{l}{\textbf{Part I: Diagnostic}} \\
All-peaks SS positive & Standard extraction & SS $= +11.7$ & --- & Fragile \\
SS robust to freq range & [1,\,85] re-extraction & $+11.7 \to -67.4$ & --- & \textbf{No} \\
SS robust to fitting window & $\phisym$-octave $\pm$ offset & $+119.3 \to -84.9$ & --- & \textbf{No} \\
Amplitude weighting rescues & 4 transforms & All attenuate & --- & \textbf{No} \\
Attractor age erosion & Pearson $r$ & $r = -0.195$ & .008 & Fragile \\
Age robust to $\fzero$ & $\fzero = 8.50$ & $r = -0.007$ & .920 & \textbf{No} \\
Age robust to IAF correction & Per-subject $\fzero$ & $r = +0.167$ (rev.) & .025 & \textbf{No} \\
$\phisym$ best (anchor-free) & Pairwise ratios & Rank 4th & --- & \textbf{No} \\

\midrule
\multicolumn{5}{l}{\textbf{Part II: What Survives}} \\
Individual alignment (OT-EO) & One-sample $t$ & $t = -5.68$, $d = 0.40$ & $< 10^{-7}$ & \textbf{Yes} \\
Individual alignment (OT-EC) & One-sample $t$ & $d = 0.75$ & $< 10^{-21}$ & \textbf{Yes} \\
$\phisym$ cross-base & Degree-3 $\bar{d}/\bar{d}_{\text{null}}$ & 0.917 (L, rank 1), 0.914 (D, rank 1), 0.901 (E, rank 1) & $< 10^{-3}$ & \textbf{Yes} \\
Cross-dataset replication & Rank 1/9 all 3 datasets & $N = 919$ total; 2nd-place base varies & $< 0.001$ & \textbf{Yes} \\
$\bar{d}$ three-dataset convergence & Mean (4-pos) & 0.0689--0.0693 & --- & \textbf{Yes} \\
Age-invariant (LEMON) & Pearson $r$ & $r = +0.054$ & 0.45 & \textbf{Yes} (null) \\
Age-invariant (Dortmund) & Pearson $r$ ($N = 601$) & $r = -0.004$ & 0.92 & \textbf{Yes} (null) \\
Cognitive prediction & 100+ tests, FDR & All null ($|r|_{\min} = 0.20$, 80\% power) & --- & \textbf{No} \\
Fatigue-invariant (Dort.) & Degree-3 pre vs post & 0.914 $\to$ 0.909 & --- & \textbf{Yes} (null) \\
Position-count decomposition & Deg 2/3/4 & $d$: 0.40 $\to$ 0.28 $\to$ 0.07 & --- & \textbf{Yes} \\
Signal at 3 aggregate positions & Per-position enrich. & Noble$_1$, attractor, boundary (aggregate; per-band may differ) & --- & \textbf{Yes} \\
EC $\times$ OT multiplicative & OT-EC vs OT-EO & $d = 0.75$ vs $0.40$ & $10^{-21}$ vs $10^{-8}$ & \textbf{Yes} \\
EC theta $\to$ boundary & Paired $t$ & $t = -3.27$ (L); $d = -0.39$ (D) & 0.001; $< 10^{-19}$ & \textbf{Yes} \\
Sign-flip disambiguation & $\fzero$ vs IAF/2 (Dort.) & $d = -0.385$ vs $+0.351$ & $< 10^{-16}$ & \textbf{Yes} \\
Cross-dataset $\theta$ convergence & OT-EO medians & $\Delta = 0.009$~Hz & --- & \textbf{Yes} \\
Target $= \fzero$ (not IAF) & Partial $r$, Williams & $r = -0.213$; $t = 4.24$ & .004; $< .0001$ & \textbf{Yes} \\
4-band lattice optimum at $\fzero$ & $\fzero$ sweep & Optimum $= 7.88$~Hz & --- & \textbf{Yes} \\
Individual pre-post $r$ (Dort.) & Pearson $r$ & $r = 0.131$ & --- & Weak \\
5-year ICC (Dort.\ $N = 208$) & ICC(2,1) & $-0.25$ to $-0.36$ & --- & \textbf{Not a trait} \\
5-year pop.\ invariance & Group $\bar{d}$ ses-2 & 0.030--0.033 & --- & \textbf{Yes} \\
5-year $\phisym$ rank (ses-2) & Cross-base rank & 1st--6th (varies by condition) & mixed & Variable \\
5-year fatigue invar.\ (ses-2) & Pre vs post $\Delta$ & $< 0.001$ & --- & \textbf{Yes} (null) \\
5-year age-mod.\ stability & $|\Delta\bar{d}|$ vs age & $r = -0.03$ to $-0.08$ & $> 0.27$ & \textbf{Yes} (null) \\

\end{longtable}

% =====================================================================
% FIGURES
% =====================================================================

\clearpage
\section{Figures}

\begin{figure}[H]
\centering
\includegraphics[width=\textwidth]{fig1_lattice_alignment.png}
\caption{\textbf{The $\phisym$-lattice and individual alignment.} (a)~Six degree-3 lattice positions on the unit circle: boundary ($u = 0$), degree-3 nobles ($0.236$, $0.764$), degree-1 nobles ($0.382$, $0.618$), and attractor ($0.5$). At degree 2, $\phisym$'s self-similarity collapses positions to 4; degree 3 restores 6 distinct positions. (b)~Per-subject mean lattice distance ($\bar{d}$) distribution vs analytical null expectation (degree 2: $t = -5.68$, $d = 0.40$, $p < 10^{-7}$). (c)~Per-band distance distributions showing theta and alpha drive the alignment signal. (d)~Cross-base comparison using degree-3 symmetric positions with coverage normalization ($\bar{d}/\bar{d}_{\text{null}}$): $\phisym$ ranks 1st in all three datasets (0.901--0.917; see Figure~\ref{fig:three_dataset}). Dominant peaks fall 8--10\% closer to $\phisym$-lattice positions than chance.}
\label{fig:alignment}
\end{figure}

\begin{figure}[H]
\centering
\includegraphics[width=\textwidth]{fig2_diagnostic_cascade.png}
\caption{\textbf{The diagnostic cascade: all-peaks metrics are fragile.} (a)~Structural score collapses when FOOOF frequency range extends to [1,\,85]~Hz, in both LEMON and EEGMMIDB datasets. (b)~Per-$\phisym$-octave fitting produces enrichment that tracks fitting window edges: shifting the band boundary by 0.5 $\phisym$-octave inverts all enrichment signs. (c)~Overlap-trim improves stability but attractor enrichment still flips sign across offsets. (d)~\texttt{max\_n\_peaks}~$= 40$ eliminates attractor age trend significance.}
\label{fig:diagnostic}
\end{figure}

\begin{figure}[H]
\centering
\includegraphics[width=\textwidth]{fig3_age_iaf_positioning.png}
\caption{\textbf{The age effect is IAF positioning.} (a)~Band decomposition shows alpha and beta drive the attractor age correlation; delta+theta is null ($p = 0.10$), confirming the effect is IAF slowing in lattice coordinates. (b)~Attractor enrichment vs age at two anchor frequencies: $\fzero = 7.83$ shows a correlation ($r = 0.31$) that vanishes at $\fzero = 8.50$ ($r \approx 0$). (c)~Per-subject $\fzero$ correction reverses the sign of the age effect. (d)~Mechanism: young and elderly IAF map to different lattice coordinates at fixed $\fzero$.}
\label{fig:age_iaf}
\end{figure}

\begin{figure}[H]
\centering
\includegraphics[width=\textwidth]{fig4_age_invariance_cognitive_null.png}
\caption{\textbf{Age-invariance and cognitive silence.} (a)~Mean lattice distance ($\bar{d}$) vs age: flat relationship ($r = 0.054$, $p = 0.45$); note that the per-subject mean $\bar{d} = 0.069$ reported in panel~(a) is the average of 202 individual $\bar{d}$ values, each computed from that subject's four dominant peaks. (b)~Population median peaks mapped to $\phisym$-lattice coordinates at $\fzero = 7.83$~Hz, showing theta and alpha near noble$_1$; the displayed $\bar{d} = 0.026$ is instead the lattice distance of the four population-median frequencies (one per band), which is lower because inter-subject variability cancels. (c)~Age-partialed correlations between $\bar{d}$ and nine cognitive measures: none reach significance. (d)~Cognitive null across three independent metrics (all within LEMON): zero FDR survivors across all-peaks, within-band $z$-scores, and dominant-peak distance.}
\label{fig:cognitive_null}
\end{figure}

\begin{figure}[H]
\centering
\includegraphics[width=\textwidth]{fig5_theta_migration.png}
\caption{\textbf{Theta position migration under eyes-closed conditions.} (a)~Theta position transition matrix (EO $\to$ EC, matched extraction, $N = 182$): boundary is 92\% stable (absorbing state), while 61\% of noble$_1$ subjects migrate to boundary. (b)~Theta frequency shift relative to $\fzero$: convergence from both directions (below-$\fzero$ subjects shift up, above-$\fzero$ subjects shift down). (c)~Position stability by starting position: boundary is maximally stable (92\% vs 32\% for noble$_1$). (d)~Power changes at positions that remain stable across conditions.}
\label{fig:theta_migration}
\end{figure}

\begin{figure}[H]
\centering
\includegraphics[width=\textwidth]{fig6_position_phenotyping.png}
\caption{\textbf{Position phenotyping.} (a)~Theta-alpha position-pair heatmap: 15 combinations observed, with boundary-noble$_1$ (19\%) and boundary-noble$_2$ (17\%) most common. (b)~Spectral profiles by theta position group: boundary subjects are stronger oscillators (higher theta and alpha power). (c)~No cognitive differences between boundary and noble groups (0/9 survive FDR).}
\label{fig:phenotyping}
\end{figure}

\begin{figure}[H]
\centering
\includegraphics[width=\textwidth]{fig7_theta_target_disambiguation.png}
\caption{\textbf{Disambiguating the EC theta convergence target: fixed $\fzero$ vs IAF-derived subharmonics.} (a)~Model-free covariance: EC$_\theta$ and EO$_\theta$ covary negatively with IAF; partial correlation controlling for EO$_\theta$ yields $r(\text{EC}_\theta, \text{IAF} \mid \text{EO}_\theta) = -0.213$ ($p = 0.004$), ruling out IAF-subharmonic predictions (which require a positive relationship). Dashed red: $\fzero = 7.83$~Hz; dotted purple: IAF/$\sqrt{\phisym}$. (b)~Prediction error distributions $|\text{EC}_\theta - \text{target}|$ comparing $\fzero = 7.83$~Hz to five IAF-derived targets; paired Wilcoxon tests confirm $\fzero$ is significantly closer (all $p < 10^{-12}$; Cliff's $\delta = 0.66$, 83\% of subjects). (c)~IAF-tercile stratification: convergence toward $\fzero$ is consistent across IAF strata ($r = -0.011$, $p = 0.88$), whereas IAF/$\sqrt{\phisym}$ predictions (purple diamonds: 7.45, 7.92, 8.42~Hz) shift with IAF by construction. Actual EC theta landing zones (filled circles) do not track them. (d)~Convergence-direction regression $\Delta\theta \sim \beta(\text{target} - \theta_{\text{EO}})$: $\fzero$ achieves $R^2 = 0.215$ vs 0.129 for IAF/$\sqrt{\phisym}$; dependent-correlation Williams test confirms superior prediction under the fixed-target model ($t = 4.24$, $p < 0.0001$).}
\label{fig:disambiguation}
\end{figure}

\begin{figure}[H]
\centering
\includegraphics[width=\textwidth]{fig8_three_dataset_replication.png}
\caption{\textbf{Three-dataset replication: $\phisym$ ranks 1st in all three datasets.} Coverage-normalized alignment ($\bar{d}/\bar{d}_{\text{null}}$) at degree-3 symmetric noble positions for nine logarithmic bases across LEMON ($N = 202$, EO), Dortmund ($N = 608$, EC-pre), and EEGMMIDB ($N = 109$). Values below 1.0 (dashed line) indicate dominant peaks closer to lattice positions than chance. $\phisym$ ranks 1st in all three datasets; the second-place base varies (1.7, 3/2, $\sqrt{2}$). Error bars: 95\% bootstrap CI.}
\label{fig:three_dataset}
\end{figure}

\begin{figure}[H]
\centering
\includegraphics[width=\textwidth]{fig9_measurement_sensitivity.png}
\caption{\textbf{Measurement sensitivity vs architectural invariance.} (a)~Within-Dortmund degree-2 condition gradient: Cohen's $d$ for individual-level alignment varies monotonically with condition (EO-pre: $d = 0.04$, ns; EO-post: $d = 0.21$; EC-post: $d = 0.27$; EC-pre: $d = 0.39$). Same subjects, same equipment---the architecture is constant but individual-level detection requires sufficient spectral SNR. (b)~Three-dataset degree-3 $\bar{d}/\bar{d}_{\text{null}}$ for $\phisym$: EC values cluster near 0.91 and EO values near 0.97. The coverage-normalized architectural metric is relatively stable across conditions, with $\phisym$ ranking 1st in all conditions except EO-pre (rank 2, behind 1.8 by 0.002).}
\label{fig:measurement}
\end{figure}

\begin{figure}[H]
\centering
\includegraphics[width=\textwidth]{fig10_disambiguation_signflip.png}
\caption{\textbf{Dortmund sign-flip: EC theta converges on $\fzero$ and diverges from IAF/2.} (a)~Distribution of $|\theta_{\text{EC}} - \fzero| - |\theta_{\text{EO}} - \fzero|$: negative values indicate EC theta closer to $\fzero$ ($d = -0.385$, $p < 10^{-19}$). (b)~Distribution of $|\theta_{\text{EC}} - \text{IAF}/2| - |\theta_{\text{EO}} - \text{IAF}/2|$: positive values indicate EC theta farther from IAF/2 ($d = +0.351$, $p < 10^{-16}$). (c)~The sign flip summarized: EC theta moves toward $\fzero$ and away from IAF/2 simultaneously, excluding all IAF-derived convergence models.}
\label{fig:signflip}
\end{figure}

\begin{figure}[H]
\centering
\includegraphics[width=\textwidth]{fig11_crossdataset_theta.png}
\caption{\textbf{Cross-dataset EO theta convergence under overlap-trim extraction.} Under standard extraction, LEMON theta starts at 7.64~Hz (above $\fzero$) and Dortmund theta starts at 5.00~Hz (below $\fzero$). Under overlap-trim extraction, both converge to $\sim$6.56~Hz ($\Delta = 0.009$~Hz), at the noble$_1$ lattice position ($u \approx 0.633$). Gold band: convergence zone. Dashed grey: $\fzero = 7.83$~Hz.}
\label{fig:theta_convergence}
\end{figure}

\begin{figure}[H]
\centering
\includegraphics[width=\textwidth]{fig12_position_count.png}
\caption{\textbf{Position-count decomposition validates the dominant-peak signal.} Cohen's $d$ (alignment vs null) attenuates monotonically from degree 2 (4 positions, $d \approx 0.38$--$0.40$) through degree 3 (6 positions, $d \approx 0.28$--$0.34$) to degree 4 (8 positions, $d \approx 0$--$0.08$) across all three datasets. Adding positions that carry no real signal systematically dilutes the metric, confirming that the alignment is concentrated at the four core lattice positions.}
\label{fig:poscount}
\end{figure}

% =====================================================================
% SUPPLEMENTARY FIGURES
% =====================================================================

\clearpage
\section*{Supplementary Figures}
\addcontentsline{toc}{section}{Supplementary Figures}

\begin{figure}[H]
\centering
\includegraphics[width=\textwidth]{figS1_amplitude_weighting.png}
\caption{\textbf{Figure S1: Amplitude weighting attenuates the signal.} Three weighting schemes (unweighted, rank, z-score) applied to attractor enrichment vs age. All weight transforms attenuate the age correlation relative to the unweighted baseline. Per-subject $\fzero$ correction (IAF anchor) reverses the sign entirely ($r = +0.167$), confirming the age effect operates through IAF-to-lattice mapping, not through peak amplitude.}
\label{fig:s1}
\end{figure}

\begin{figure}[H]
\centering
\includegraphics[width=\textwidth]{figS2_base_specificity.png}
\caption{\textbf{Figure S2: Base specificity in anchor-free analysis.} (a)~RWT $\sim$ $z_{\text{alpha}}$ correlation appears at multiple bases: 3/2 ($r = 0.309$) beats $\phisym$ ($r = 0.227$), and 5 of 9 bases survive FDR. The cognitive effect is generic alpha-band spectral organization, not $\phisym$-specific. (b)~Summary: dominant-peak mean $\bar{d}$ is the only metric where $\phisym$ achieves rank 1.}
\label{fig:s2}
\end{figure}

\begin{figure}[H]
\centering
\includegraphics[width=\textwidth]{figS3_extraction_sensitivity.png}
\caption{\textbf{Figure S3: Extraction method sensitivity.} (a)~Position agreement by band between standard and overlap-trim extraction on the same EO data. (b)~Summary comparison: standard vs overlap-trim achieves 66\% agreement, worse than the 74\% EO--EC agreement with matched extraction. One-third of subjects change theta position depending on FOOOF fitting approach.}
\label{fig:s3}
\end{figure}

% =====================================================================
% APPENDICES
% =====================================================================

\clearpage
\appendix

\section{Relationship to Papers 1 and 2}
\label{app:papers12}

This is the third paper in a series. Paper~1 established the $\phisym^n$ lattice framework through four complementary studies---Schumann Ignition Event (SIE) discovery, single-channel spectral architecture, multi-channel spatial coherence, and independent replication---demonstrating that human neural oscillations organise along $f(n) = \fzero \times \phisym^n$ with convergent $\fzero$ estimates near 7.6~Hz. Paper~2 provided extended structural specificity analysis using the EEGMMIDB dataset (109 subjects, 64 channels), including base comparison across nine exponential lattices and anchor-frequency optimisation. This paper applies the framework to the LEMON dataset ($N = 202$, 62 channels) with diagnostic cascade, then replicates the surviving findings in the Dortmund Vital Study ($N = 608$, 64 channels, ages 20--70) and revises several claims from Papers~1 and~2.

\subsection{Claims retained}

\begin{itemize}
    \item \textbf{Irrational $>$ rational bases} (Papers~1--2): Irrational bases ($\phisym$, $e$, $\pi$, $\sqrt{2}$) outperform rational bases (3/2, 2) on aggregate enrichment ($d = 1.2$--$1.6$). Replicated in LEMON. This validates the logarithmic-scaling hypothesis independent of the choice among irrational bases.
    \item \textbf{Non-exponential grids null} (Paper~2): Linear harmonic grids produce null structural scores (mean SS $\approx 0$). This confirms that the enrichment signal (such as it is) requires logarithmic spacing.
    \item \textbf{State modulation is not $\phisym$-specific} (Paper~2): The gamma reallocation effect persists across all anchor frequencies, all bases except $e$, and random positions ($p_{\text{emp}} = 0.44$). Replicated in LEMON EC/EO comparison.
\end{itemize}

\subsection{Claims revised}

\begin{itemize}
    \item \textbf{``$\phisym$ is the global optimum across all bases and anchor frequencies''} (Paper~2, SS $= 45.6$). \textit{Correction:} This all-peaks SS depended on the FOOOF frequency range [1,\,50]~Hz, which truncates the gamma $\phisym$-octave and biases the within-band shuffle null. Extending to [1,\,75]~Hz collapses SS to $-23.6$ (Section~3.1). The same collapse occurs in LEMON ([1,\,85]~Hz: $+11.7 \to -67.4$). Paper~2's revision retains SS $= 45.6$ as conditional on anchor range ($\fzero = 7.6$--$8.5$~Hz) and labels it as an aggregate cross-band figure. $\phisym$ remains the best base for \textit{dominant-peak} alignment in a bounded anchor window, but the unrestricted all-peaks optimality claim is retracted.
    \item \textbf{``The optimal anchor is $\fzero = 8.5$~Hz''} (Paper~2, all-peaks joint grid search). \textit{Correction:} This all-peaks optimum was driven by the gamma truncation artifact. Under the dominant-peak analysis introduced in this paper, the lattice optimum is $\fzero = 7.88$~Hz (Section~\ref{sec:anchor_confirmation}), consistent with the pre-specified Schumann fundamental ($\fzero = 7.83$). Paper~2 retains $\fzero = 8.5$~Hz as the all-peaks optimum within its conditional framing; the dominant-peak metric supersedes this for anchor determination.
    \item \textbf{All-peaks metrics dependent on FOOOF parameters.} Any claim from Papers~1 or~2 that depends on specific \texttt{max\_n\_peaks}, FOOOF frequency range, or post-hoc anchor optimisation should not be cited without the diagnostic qualifications presented here. These parameters control the sign and magnitude of all-peaks metrics (Section~3.1--3.2). Consistent with this finding, the revisions to Papers~1 and~2 now label aggregate enrichment values (boundary $-18\%$, Noble$_1$ $+39\%$ in Paper~1; boundary $0.41\times$, Noble$_1$ $1.10\times$ in Paper~2) as ``aggregate cross-band'' figures, acknowledge per-band heterogeneity in structural scores, and identify per-band enrichment normalization as a priority for future validation.
\end{itemize}

\subsection{New findings from this paper}

\begin{itemize}
    \item \textbf{Three-dataset dominant-peak alignment} ($N = 919$): Per-subject $\bar{d}$ converges to $0.069 \pm 0.0004$ across LEMON, Dortmund, and EEGMMIDB. Under degree-3 symmetric coverage normalization, $\phisym$ ranks 1st in all three datasets ($\bar{d}/\bar{d}_{\text{null}} = 0.901$--$0.917$). Age-invariant ($r = -0.004$, Dortmund continuous 20--70), fatigue-invariant (0.914 $\to$ 0.909), cognitively silent.
    \item \textbf{Definitive $\fzero$ disambiguation via sign flip}: Dortmund's geometry places $\fzero$ and IAF/2 in opposite directions from theta. EC theta converges on $\fzero$ ($d = -0.385$) while diverging from IAF/2 ($d = +0.351$), excluding all IAF-derived models.
    \item \textbf{Cross-dataset theta convergence}: LEMON (7.64~Hz) and Dortmund (5.00~Hz) EO theta converge to 6.56~Hz ($\Delta = 0.009$~Hz) under overlap-trim extraction, from opposite directions.
    \item \textbf{Position-count decomposition}: $d \approx 0.40$ (4 pos) $\to$ 0.28 (6 pos) $\to$ 0.07 (8 pos) across all three datasets---monotonic attenuation validates real signal.
    \item \textbf{Crystal-not-fingerprint}: Population architecture is invariant (degree-3 $\bar{d}/\bar{d}_{\text{null}} \approx 0.91$), but individuals reshuffle (pre-post $r = 0.131$). The lattice describes allowed sites, not occupied ones.
    \item \textbf{Measurement sensitivity vs architectural invariance}: Degree-2 Cohen's $d$ varies with spectral SNR (within-Dortmund gradient: 0.037--0.391). Degree-3 $\bar{d}/\bar{d}_{\text{null}}$ is condition- and dataset-invariant ($\sim$0.91).
    \item \textbf{EC strengthens four-band alignment}: All four band modes tighten under EC ($d = 0.75$, $p = 10^{-21}$), multiplicative with overlap-trim extraction.
    \item \textbf{Comprehensive extraction diagnostic}: All-peaks enrichment is dominated by FOOOF parameters. The attractor age erosion ($r = -0.195$) is fully explained by IAF positioning.
\end{itemize}

\section{Detailed FOOOF sensitivity analysis}

\subsection{\texttt{max\_n\_peaks} sensitivity}

At \texttt{max\_n\_peaks}~$= 20$: 215K peaks, attractor age $r = -0.195$ ($p = 0.008$). At \texttt{max\_n\_peaks}~$= 40$: 405K peaks, boundary enrichment inverts to $+19.7\%$, attractor age $r = -0.110$ ($p = 0.138$). Approximately 54\% of channels hit the 20-peak ceiling at LEMON's 0.122~Hz resolution. Spearman rank correlation of per-subject compliance: $\rho(\text{max20}, \text{max40}) = 0.838$.

\subsection{Raw vs preprocessed extraction}

\begin{table}[H]
\centering
\begin{tabular}{lcc}
\toprule
Property & Preprocessed & Raw \\
\midrule
Aperiodic exponent & 2.45 & 1.20 \\
Peak freq range & 4.5--71.0~Hz & 1.1--85.0~Hz \\
SS (global, [1,\,85]) & $-67.4$ & $+68.5$ \\
$N$ subjects & 202 & 208 \\
\bottomrule
\end{tabular}
\end{table}

ICA removes broadband activity, steepening the aperiodic slope. Raw data preserves the full spectral range but includes muscle and eye artifacts. The sign reversal in SS between raw ($+68.5$) and preprocessed ($-67.4$) further demonstrates metric fragility.

\section{Disambiguation robustness analyses}
\label{app:disambiguation_robustness}

\subsection{Best-constant sensitivity to theta estimator}

The best-fit constant (the single frequency minimizing MSE to EC theta) varies with estimation method: 7.38~Hz (dominant FOOOF peak, default), 7.44 (posterior-only dominant peak), 7.15 (frontal-only dominant peak), and 6.43 (power-weighted center-of-gravity across all theta peaks). Among dominant-peak estimators---the class relevant to the lattice analysis---the range is 7.15--7.44~Hz (spread $= 0.29$~Hz). The CoG outlier (6.43~Hz) reflects averaging toward band center rather than a different convergence target. All estimators place the best-fit constant closer to $\fzero = 7.83$ than to any IAF-derived prediction (IAF/$\sqrt{\phisym} \approx 7.93$, IAF$\times$3/4 $\approx 7.57$).

\subsection{Age stratification}

The $\fzero$ advantage holds within both age groups:

\begin{itemize}
    \item \textbf{Young} ($N = 123$): $\fzero$ closer for 83\% of subjects, $R^2 = 0.214$ vs 0.128 for IAF/$\sqrt{\phisym}$, partial $r(\text{EC}_\theta, \text{IAF} \mid \text{EO}_\theta) = -0.225$ ($p = 0.013$).
    \item \textbf{Elderly} ($N = 59$): $\fzero$ closer for 83\%, $R^2 = 0.209$ vs 0.121, partial $r = -0.207$ ($p = 0.12$---non-significant due to reduced power but same direction and magnitude).
\end{itemize}

With age as a covariate in the full sample: convergence-direction effect unchanged (partial $r(\Delta\theta, \fzero\text{-gap} \mid \text{age}) = 0.464$, $p < 10^{-10}$); IAF-invariance survives ($r = -0.014$, $p = 0.85$). The disambiguation conclusion is not confounded by age.

\section{Per-$\phisym$-octave fitting edge artifacts}
\label{app:peroctave}

Fitting FOOOF separately within each $\phisym$-octave band (to match analysis windows with extraction) produces enrichment that tracks fitting window edges, not lattice positions:
\begin{itemize}
    \item Primary (edges at lattice boundaries): SS $= +119.3$, $E_{\text{boundary}} = -60.8\%$, $E_{\text{attractor}} = +31.6\%$
    \item Offset by $0.5$ $\phisym$-octave: SS $= -84.9$, $E_{\text{boundary}} = +14.8\%$, $E_{\text{attractor}} = -48.2\%$
\end{itemize}
The boundary depletion completely inverts ($-60.8\% \to +14.8\%$) when band edges shift by half a $\phisym$-octave. Enrichment is a property of the fitting window, not the EEG.

% =====================================================================
% REFERENCES
% =====================================================================

\printbibliography

\end{document}
